\documentclass[12pt]{article}
\usepackage[ukrainian,shorthands=off]{babel}   % shorthands=off: символ " не активний (потрібно для міток "..." у TikZ всередині \zadtask)
\usepackage{fontspec}
% --- НАЛАШТУВАННЯ ШРИФТІВ ---
% Шрифти Schoolbook-*.otf лежать у корені проєкту. Шлях підбирається автоматично:
% ../ (компіляція з папки теми), ./ (корінь проєкту), ../../ (підпапки теми 28); інакше -- TeX Gyre Schola.
\newcommand{\nmtSchoolbook}[1]{\setmainfont{Schoolbook-Regular.otf}[
    Path = #1,
    BoldFont = Schoolbook-Bold.otf,
    ItalicFont = Schoolbook-Italic.otf,
    BoldItalicFont = Schoolbook-BoldItalic.otf
]}
\IfFileExists{../Schoolbook-Regular.otf}{\nmtSchoolbook{../}}{%
 \IfFileExists{./Schoolbook-Regular.otf}{\nmtSchoolbook{./}}{%
  \IfFileExists{../../Schoolbook-Regular.otf}{\nmtSchoolbook{../../}}{%
   \setmainfont{texgyreschola}[Extension=.otf, UprightFont=*-regular, BoldFont=*-bold, ItalicFont=*-italic, BoldItalicFont=*-bolditalic]}}}
\usepackage[italic]{mathastext}
\usepackage[a4paper,margin=1.8cm,bottom=2cm,top=2cm]{geometry}
\usepackage{amsmath,amssymb}
\usepackage{tikz}
\usepackage{pgfplots}
\pgfplotsset{compat=1.18}
\usepgfplotslibrary{fillbetween}
\usetikzlibrary{calc,patterns,angles,quotes,intersections,babel,3d,shapes.symbols,shapes.geometric,decorations.pathreplacing,shadings,arrows.meta,hobby}
% ===================== ДИЗАЙН-СИСТЕМА РИСУНКІВ (scripts/ПРАВИЛА_ОФОРМЛЕННЯ.md, scripts/дизайн_завдань.py) =====================
\definecolor{nmtInk}{HTML}{1A1A1A}
\definecolor{nmtGrid}{HTML}{3A3A3A}      % сітка, як у НМТ 2026: темна й тонка (0,35 pt), клітинка 0,5 см
\definecolor{nmtBlue}{HTML}{1F6FB2}\definecolor{nmtGreen}{HTML}{2E8B57}\definecolor{nmtRed}{HTML}{C0392B}
\definecolor{nmtOrange}{HTML}{D9822B}\definecolor{nmtViolet}{HTML}{7D4CA8}
\definecolor{nmtFillA}{HTML}{D3E5F5}\definecolor{nmtFillB}{HTML}{D8EDD5}\definecolor{nmtFillC}{HTML}{FBE1C8}
\definecolor{nmtFillD}{HTML}{E4E4E4}\definecolor{nmtFillE}{HTML}{E7DCF2}
\tikzset{
  ultra thin/.style={line width=0.2pt}, very thin/.style={line width=0.3pt}, thin/.style={line width=0.45pt},
  semithick/.style={line width=0.6pt}, thick/.style={line width=0.75pt}, very thick/.style={line width=1.1pt},
  ultra thick/.style={line width=1.5pt},
  >={Stealth[length=2.3mm, width=1.7mm]},
  every picture/.append style={line join=round, line cap=round},
  nmt dot/.style={circle, fill=nmtInk, draw=none, inner sep=0pt, minimum size=3.6pt},
  nmt open dot/.style={circle, fill=white, draw=nmtInk, line width=0.6pt, inner sep=0pt, minimum size=4.2pt},
  nmt grid/.style={draw=nmtGrid, line width=0.35pt},
  nmt curve/.style={line width=1.1pt},
}
% вирази: ≤ і ≥ -- похилі, як у НМТ (десяткову кому без пробілу вже дає mathastext)
\let\leq\leqslant \let\le\leqslant \let\geq\geqslant \let\ge\geqslant
\pgfplotsset{every axis/.append style={axis line style={line width=0.6pt}, tick style={line width=0.45pt},
  tick label style={font=\footnotesize}, label style={font=\small}, grid style={nmtGrid, line width=0.35pt}}}
% ===================== кінець дизайн-системи =====================
% підпис кута біля дуги: \nmtAngleLabel{A}{B}{C}{радіус дуги, см}{підпис} -- кут AB C проти годинникової стрілки, як \pic{angle=A--B--C}.
% Підпис ставиться на бісектрисі, найближчим краєм на 0,08 см від дуги (радіус дуги масштабується разом з рисунком, підпис -- ні),
% тож не відлітає від кута, хоч який масштаб рисунка й хоч яка довжина підпису
\newcommand{\nmtAngleLabel}[5]{%
  \path let \p1=(#2), \p2=(#1), \p3=(#3), \p4=(0,0), \p5=(1,0),
    \n1={atan2(\y2-\y1,\x2-\x1)}, \n2={atan2(\y3-\y1,\x3-\x1)},
    \n3={mod(\n2-\n1+720,360)}, \n4={\n1+\n3/2}, \n5={\n4+180},
    \n6={veclen(\x5-\x4,\y5-\y4)/28.45274}, \n7={(#4*\n6+0.08)*28.45274}
  in (#2) node[anchor=\n5, inner sep=1.2pt, shift={(\n4:\n7)}] {#5};}

\usepackage{xcolor,array,fancyhdr,enumitem,multicol}
\usepackage{colortbl,multirow,diagbox,needspace}
\usepackage{tcolorbox}
\tcbuselibrary{skins,breakable}
\usepackage[hidelinks]{hyperref}
% водяний знак; щоб зібрати без нього (версія для вчителів), перед \documentclass пишуть \def\nmtnowatermark{}
\ifdefined\nmtnowatermark\else
\usepackage[firstpage=false, text=@pvtr2525, scale=0.6, color=gray!12]{draftwatermark}
\fi

% --- АВТОСТИСКАННЯ: рисунок/сітка, ширша за колонку, масштабується до ширини колонки ---
\newsavebox{\nmtfitbox}
\newenvironment{nmtfit}{\begin{lrbox}{\nmtfitbox}}{\end{lrbox}%
  \ifdim\wd\nmtfitbox>\linewidth \resizebox{\linewidth}{!}{\usebox{\nmtfitbox}}\else\usebox{\nmtfitbox}\fi}
\let\nmtIncludegraphicsOrig\includegraphics
\renewcommand{\includegraphics}[2][]{\begin{nmtfit}\nmtIncludegraphicsOrig[#1]{#2}\end{nmtfit}}


\definecolor{mainGreen}{RGB}{34, 120, 64}
\definecolor{yearOrange}{RGB}{220, 100, 30}
\definecolor{yearcolor}{RGB}{220, 100, 30}
\definecolor{titleBg}{RGB}{225, 240, 220}
\definecolor{instrBg}{RGB}{255, 240, 220}
\definecolor{instrBorder}{RGB}{220, 150, 80}
\definecolor{headerblue}{RGB}{0, 102, 204}   % використовується в рисунках старої бази

\setlength{\parindent}{0pt}
\emergencystretch=1.5em

\pagestyle{fancy}
\fancyhf{}
\renewcommand{\headrulewidth}{0.4pt}
\renewcommand{\footrulewidth}{0.4pt}
\fancyhead[C]{\small\color{gray!80} Функції: завдання на відповідність}
\fancyhead[R]{\small\color{mainGreen}\textbf{\href{https://t.me/pvtr2525}{@pvtr2525}}}
\fancyfoot[L]{\small\color{gray!80}\href{https://t.me/pvtr2525}{ТГ: @pvtr2525}}
\fancyfoot[C]{\small\color{gray!80}\thepage}
\fancyfoot[R]{\small\color{gray!80}НМТ 2022--2026}

% --- ТАБЛИЦІ ВІДПОВІДЕЙ ---
% ширина клітинки = (ширина рядка - відступи) / 5: на всю сторінку це ~3 см, у minipage поруч із рисунком таблиця стискається;
% п'ять колонок |m{w}| мають 10 відступів \tabcolsep і 6 вертикальних ліній -- тоді таблиця рівно на \linewidth
\newlength{\nmtcellw}
\newcommand{\nmtSetCell}{\setlength{\nmtcellw}{\dimexpr(\linewidth-10\tabcolsep-6\arrayrulewidth)/5\relax}}
\newcommand{\answerTable}[5]{%
\par\nopagebreak\vspace{0.15cm}
\noindent\nmtSetCell
\begin{tabular}{|*{5}{>{\centering\arraybackslash}m{\nmtcellw}|}}
\hline
\rule[-0.15cm]{0pt}{0.6cm}\textbf{А} & \rule[-0.15cm]{0pt}{0.6cm}\textbf{Б} & \rule[-0.15cm]{0pt}{0.6cm}\textbf{В} & \rule[-0.15cm]{0pt}{0.6cm}\textbf{Г} & \rule[-0.15cm]{0pt}{0.6cm}\textbf{Д} \\
\hline
\rule[-0.2cm]{0pt}{0.75cm}#1 & \rule[-0.2cm]{0pt}{0.75cm}#2 & \rule[-0.2cm]{0pt}{0.75cm}#3 & \rule[-0.2cm]{0pt}{0.75cm}#4 & \rule[-0.2cm]{0pt}{0.75cm}#5 \\
\hline
\end{tabular}
\par\vspace{0.25cm}
}
\newcommand{\answerTableTall}[5]{%
\par\nopagebreak\vspace{0.15cm}
\noindent\nmtSetCell
\begin{tabular}{|*{5}{>{\centering\arraybackslash}m{\nmtcellw}|}}
\hline
\rule[-0.2cm]{0pt}{0.7cm}\textbf{А} & \rule[-0.2cm]{0pt}{0.7cm}\textbf{Б} & \rule[-0.2cm]{0pt}{0.7cm}\textbf{В} & \rule[-0.2cm]{0pt}{0.7cm}\textbf{Г} & \rule[-0.2cm]{0pt}{0.7cm}\textbf{Д} \\
\hline
\rule[-0.45cm]{0pt}{1.2cm}#1 & \rule[-0.45cm]{0pt}{1.2cm}#2 & \rule[-0.45cm]{0pt}{1.2cm}#3 & \rule[-0.45cm]{0pt}{1.2cm}#4 & \rule[-0.45cm]{0pt}{1.2cm}#5 \\
\hline
\end{tabular}
\par\vspace{0.25cm}
}
\newcommand{\instructionBox}[1]{%
\par\vspace{0.3cm}
\begin{tcolorbox}[colback=instrBg, colframe=instrBorder, boxrule=0.6pt, arc=2pt,
    left=8pt, right=8pt, top=4pt, bottom=4pt]
\centering\small\bfseries #1
\end{tcolorbox}
\par\vspace{0.15cm}
}
\newcommand{\sectionTitle}[1]{%
\par\vspace{0.4cm}
\begin{tcolorbox}[colback=titleBg, colframe=titleBg, boxrule=0pt, arc=8pt,
    halign=center, fontupper=\Large\bfseries\color{mainGreen}]
\hyphenpenalty=10000 \exhyphenpenalty=10000 #1
\end{tcolorbox}
\par\vspace{0.3cm}
}
\newcommand{\task}[2]{\noindent\textbf{#1.}\ \ #2\par\vspace{0.2cm}}
% --- НМТ-СТИЛЬ ПОЛЕ ВІДПОВІДІ ---
\newcommand{\nmtAnswerBox}{%
\par\nopagebreak\vspace{0.2cm}\noindent
Відповідь:\ \
\begin{tabular}{@{}*{4}{|>{\centering\arraybackslash}p{0.8cm}}|@{\hspace{0.1cm}\textbf{,}\hspace{0.1cm}}*{3}{|>{\centering\arraybackslash}p{0.8cm}}|@{}}
\hline
\rule[-0.2cm]{0pt}{0.7cm} & & & & & & \\
\hline
\end{tabular}
\par\vspace{0.3cm}
}
% --- СІТКА ДЛЯ МАТЧИНГУ ---
\newsavebox{\nmtgridbox}
\newcommand{\matchingGrid}{%
    \sbox{\nmtgridbox}{\matchingGridRaw}%
    \ifdim\wd\nmtgridbox>\linewidth \resizebox{\linewidth}{!}{\usebox{\nmtgridbox}}\else\usebox{\nmtgridbox}\fi}
\newcommand{\matchingGridRaw}{%
    \begingroup
    \renewcommand{\arraystretch}{1.15}
    \setlength{\tabcolsep}{4pt}
    \footnotesize
    \begin{tabular}{r|c|c|c|c|c|}
         \multicolumn{1}{c}{} & \multicolumn{1}{c}{\textbf{А}} & \multicolumn{1}{c}{\textbf{Б}} & \multicolumn{1}{c}{\textbf{В}} & \multicolumn{1}{c}{\textbf{Г}} & \multicolumn{1}{c}{\textbf{Д}} \\ \cline{2-6}
         \textbf{1} & \phantom{X} & \phantom{X} & \phantom{X} & \phantom{X} & \phantom{X} \\ \cline{2-6}
         \textbf{2} & & & & & \\ \cline{2-6}
         \textbf{3} & & & & & \\ \cline{2-6}
    \end{tabular}
    \endgroup
}
% --- ДОДАТКОВІ МАКРОСИ ЗБІРНИКА ---
\newcommand{\taskBlock}[1]{%
\par\vspace{0.45cm}
\noindent{\color{mainGreen}\rule{\linewidth}{0.8pt}}\par\vspace{0.12cm}
\noindent{\small\bfseries\color{mainGreen}#1}\par\vspace{0.25cm}
}
\newcounter{zad}
\newcommand{\zadtask}[1]{\stepcounter{zad}\noindent\textbf{\thezad.}\ \ #1\par\nopagebreak\vspace{0.2cm}}
\newcommand{\zadnum}{\stepcounter{zad}\textbf{\thezad.}\ \ }
\newcounter{ans}
\newcommand{\ansitem}[1]{\stepcounter{ans}\noindent\textbf{\theans.}\ #1\par\vspace{0.07cm}}
% мітка року: нерозривна, притиснута праворуч; якщо не вміщається -- праворуч на новому рядку
\newcommand{\nmtyear}[1]{\unskip\nobreak\finalhyphendemerits=0 \hfill\penalty100\hskip1em\hbox{}\nobreak\hfill{\small\color{yearOrange}\mbox{\rule{0pt}{2.3ex}(НМТ~#1)}}}
% --- МАКРОС КОРОТКОГО РОЗВ'ЯЗКУ ---
\newcommand{\solution}[1]{%
\par\vspace{0.12cm}
\begin{tcolorbox}[colback=titleBg!45, colframe=mainGreen!55, boxrule=0.4pt, arc=2pt,
    left=6pt, right=6pt, top=3pt, bottom=3pt, breakable]
\footnotesize\textbf{Короткий розв'язок.}\ #1
\end{tcolorbox}
\par\vspace{0.12cm}
}
% Розв'язки й відповіді (є до завдань НМТ-2026): \showsolutionstrue -- показати, \showsolutionsfalse -- сховати
\newif\ifshowsolutions
\showsolutionsfalse
% --- ЗАГОЛОВКИ З ПУНКТАМИ КЛІКАБЕЛЬНОГО ЗМІСТУ ---
\setcounter{tocdepth}{2}
\newcommand{\chapterTitle}[1]{%
\clearpage\phantomsection\addcontentsline{toc}{section}{#1}%
\sectionTitle{#1}}
\newcommand{\typeTitle}[1]{%
\needspace{6\baselineskip}%
\phantomsection\addcontentsline{toc}{subsection}{\quad #1}%
\par\vspace{0.3cm}
\noindent{\large\bfseries\color{yearOrange}#1}\par\vspace{0.08cm}
\noindent{\color{yearOrange}\rule{\linewidth}{0.4pt}}\par\nopagebreak\vspace{0.2cm}}
\raggedbottom
% усередині одного завдання сторінка не розривається: нерозривний вертикальний проміжок
% і нерозривна межа абзаців (умова ніколи не відділяється від варіантів, рисунка чи сітки)
\newcommand{\nbvspace}[1]{\nopagebreak\vspace{#1}}
\newcommand{\nmtnobreak}{\ifvmode\penalty10000 \fi}
% --- ЄДИНИЙ МАКЕТ ЗАВДАНЬ НА ВІДПОВІДНІСТЬ ---
% мітка в колонці сталої ширини, текст -- у parbox: продовження довгого рядка
% вирівнюється під текстом, а не під міткою; відступ між пунктами однаковий скрізь
\newlength{\nmtMatchLab}\setlength{\nmtMatchLab}{1.6em}
\newlength{\nmtMatchSep}\setlength{\nmtMatchSep}{0.28cm}
\newcommand{\matchHead}[1]{\par\noindent\textit{#1}\par\nopagebreak\vspace{0.2cm}}
% шаблонний заголовок стовпця зі старого макета -- лише коли завдання не має власного \matchHead
\makeatletter\newcommand{\nmtHeadUnless}[2]{\in@{\matchHead}{#1}\ifin@\else#2\fi}\makeatother
\newcommand{\matchItem}[2]{\par\noindent\makebox[\nmtMatchLab][l]{\textbf{#1}}%
\parbox[t]{\dimexpr\linewidth-\nmtMatchLab\relax}{\raggedright #2}\par\nopagebreak\vspace{\nmtMatchSep}}
\newcommand{\ansTheme}[1]{\par\vspace{0.25cm}\noindent{\bfseries\color{mainGreen}#1}\par\vspace{0.12cm}}
\newcommand{\ansType}[1]{\par\vspace{0.1cm}\noindent{\itshape\color{yearOrange}#1}\par\vspace{0.08cm}}

% варіанти відповіді окремими рядками (довгі вирази не влазять у клітинки таблиці)
\newcommand{\answerRows}[5]{\par\nopagebreak\vspace{0.15cm}%
\matchItem{А}{#1}\matchItem{Б}{#2}\matchItem{В}{#3}\matchItem{Г}{#4}\matchItem{Д}{#5}}
\graphicspath{{./}{20. Функція та її властивості/}{21.  Графіки елементарних функцій, їх побудова графіків функцій та їх перетворення/}{28. Показникова функція. Показникові рівняння/Показникові функція і вирази/}{30. Логарифм. Логарифмічна функція/}{33. Похідна/}}

\begin{document}
\begingroup\setcounter{zad}{0}
% --- макроси з файлів сесій НМТ-2026 (для рисунків) ---
\renewcommand{\headrulewidth}{0.4pt}
\renewcommand{\footrulewidth}{0.4pt}
\definecolor{gold}{RGB}{240, 190, 40}
\definecolor{silver}{RGB}{170, 170, 175}
\definecolor{bronze}{RGB}{190, 120, 70}
\definecolor{hexcol}{RGB}{200,110,60}
\newcommand{\hexthree}[2]{%
    \draw[hexcol, line width=1.1pt]
        ($(#1,#2)+(0:30)$) -- ($(#1,#2)+(60:30)$) -- ($(#1,#2)+(120:30)$) --
        ($(#1,#2)+(180:30)$) -- ($(#1,#2)+(240:30)$) -- ($(#1,#2)+(300:30)$) -- cycle;
}

% --- сумісність зі старою базою (діє, якщо тема не має власного визначення) ---
\providecommand{\trigStyle}[1]{#1}
\providecommand{\answerTableSmall}[5]{%
\begingroup\setlength{\tabcolsep}{2pt}\small\nmtSetCell
\begin{tabular}{|*{5}{>{\centering\arraybackslash}m{\nmtcellw}|}}
\hline
\rule[-0.2cm]{0pt}{0.6cm}\textbf{А} & \textbf{Б} & \textbf{В} & \textbf{Г} & \textbf{Д} \\
\hline
\rule[-0.4cm]{0pt}{0.9cm}#1 & \rule[-0.4cm]{0pt}{0.9cm}#2 & \rule[-0.4cm]{0pt}{0.9cm}#3 & \rule[-0.4cm]{0pt}{0.9cm}#4 & \rule[-0.4cm]{0pt}{0.9cm}#5 \\
\hline
\end{tabular}\endgroup}
\providecommand{\matchingLayout}[3]{%
    \noindent
    \begin{minipage}[t]{0.36\textwidth}\vspace{0pt}\raggedright
        #1
    \end{minipage}%
    \hfill
    \begin{minipage}[t]{0.43\textwidth}\vspace{0pt}\raggedright
        #2
    \end{minipage}%
    \hfill
    \begin{minipage}[t]{0.19\textwidth}
        \vspace{0pt}
        \begin{flushright}
        #3
        \end{flushright}
    \end{minipage}}
\providecommand{\matchingLayoutBottom}[2]{%
    \noindent
    \begin{minipage}[t]{0.48\textwidth}\vspace{0pt}\raggedright
        #1
    \end{minipage}%
    \hfill
    \begin{minipage}[t]{0.48\textwidth}\vspace{0pt}\raggedright
        #2
    \end{minipage}
    \vspace{0.5cm}
    \noindent
    \matchingGrid}
\providecommand{\answerListVertical}[5]{%
    \par\nopagebreak\vspace{0.2cm}
    \begin{itemize}[itemsep=0.4cm, leftmargin=1.5cm, labelsep=0.5cm]
        \item[\textbf{А}] #1
        \item[\textbf{Б}] #2
        \item[\textbf{В}] #3
        \item[\textbf{Г}] #4
        \item[\textbf{Д}] #5
    \end{itemize}
    \vspace{0.2cm}}

\graphicspath{{./}{11. Трикутники/}}
\setcounter{zad}{0}

\chapterTitle{Функції: завдання на відповідність}

\noindent{\small Збірник містить \textbf{усі 69 завдань} НМТ 2023--2026 років на встановлення відповідності з бази (функції та їхні графіки). \textbf{Частина~1} (8) -- без тригонометрії, показникових функцій, логарифмів і похідної, \textbf{частина~2} (61) -- решта. Усередині частин завдання згруповано за типом; біля кожного вказано рік. Відповіді~--- на останній сторінці.\par\vspace{0.15cm}
За роками: \mbox{2023~--- 19}, \mbox{2024~--- 19}, \mbox{2025~--- 14}, \mbox{2026~--- 17}.}\par\vspace{0.45cm}

\begin{samepage}
\needspace{0.45\textheight}\sectionTitle{Частина 1. Без тригонометрії, показникових функцій, логарифмів і похідної}
\typeTitle{Характеристика функції --- проміжок:\\ $D(y)$, $E(y)$, проміжок монотонності}
% НМТ 2026, сесія 23.05, завдання №18
\zadtask{Установіть відповідність між характеристикою функції \mbox{(1--3)} та проміжком \mbox{(А--Д)}, якому вона належить. \nmtyear{2026}

\nmtnobreak
\nopagebreak\nbvspace{0.3cm}

\nmtnobreak
\noindent
\matchingLayout{
\matchHead{Характеристика}
\matchItem{1}{область визначення функції $y = \sqrt{x - 2}$}
\matchItem{2}{точка екстремуму функції $y = x^2 + 2x + 3$}
\matchItem{3}{абсциса точки перетину прямої $y = 6$ з графіком функції $y = x^3$}
}{
\matchHead{Проміжок}
\matchItem{А}{$(-\infty;\, -2)$}
\matchItem{Б}{$[-2;\, 0)$}
\matchItem{В}{$[0;\, 1)$}
\matchItem{Г}{$(1;\, 2)$}
\matchItem{Д}{$[2;\, +\infty)$}
}{
\matchingGrid
}
}
% Відповідь: 1~--~Д, 2~--~Б, 3~--~Г
\ifshowsolutions\solution{$1)$ ОДЗ $y=\sqrt{x-2}$: $x\ge 2$, тобто $[2;+\infty)$~(Д). $2)$ Вершина параболи $y=x^2+2x+3$: $x=-\dfrac{2}{2}=-1\in[-2;0)$~(Б). $3)$ $x^3=6 \Rightarrow x=\sqrt[3]{6}\approx 1{,}8\in(1;2)$~(Г). Відповідь: \textbf{1--Д, 2--Б, 3--Г}.}\fi
\par\penalty-20
\end{samepage}

\begin{samepage}
\needspace{0.45\textheight}\typeTitle{Кількість спільних точок графіка функції з прямою}
% НМТ 2026, сесія 19.06, завдання №17
\zadtask{До кожного початку речення \mbox{(1--3)} доберіть його закінчення \mbox{(А--Д)} так, щоб утворилося правильне твердження. \nmtyear{2026}\par
\nopagebreak\nbvspace{0.3cm}
\noindent
\matchingLayout{
\matchHead{Початок речення}
\matchItem{1}{Графіки функцій $y = 3 - x$ та $y = 1 - x$}
\matchItem{2}{Графіки функцій $y = \dfrac{1}{x}$ та $y = x^2 - 4$}
\matchItem{3}{Графік функції $y = 3$ та графік рівняння $x^2 + y^2 = 9$}
}{
\matchHead{Закінчення речення}
\matchItem{А}{не мають спільних точок.}
\matchItem{Б}{мають єдину спільну точку.}
\matchItem{В}{мають лише дві спільні точки.}
\matchItem{Г}{мають лише три спільні точки.}
\matchItem{Д}{мають чотири спільні точки.}
}{
\matchingGrid
}
}
% Відповідь: 1--А,\ \ 2--Г,\ \ 3--Б
\par\penalty-20
\end{samepage}

\begin{samepage}
% НМТ 2026, сесія 3.06, завдання №17
\noindent\zadnum До кожного початку речення \mbox{(1--3)} доберіть його закінчення \mbox{(А--Д)} так, щоб утворилося правильне твердження. \nmtyear{2026}

\nmtnobreak
\nopagebreak\nbvspace{0.3cm}
\noindent
\matchingLayout{
\matchHead{Початок речення}
\matchItem{1}{Графіки функцій $y = 2x - 3$ та $y = x + 1$}
\matchItem{2}{Графіки функцій $y = x^2$ та $y = x^2 - 2$}
\matchItem{3}{Графік функції $y = \dfrac{2}{x}$ та графік рівняння $x^2 + y^2 = 4$}
}{
\matchHead{Закінчення речення}
\matchItem{А}{не мають жодної спільної точки.}
\matchItem{Б}{мають єдину спільну точку.}
\matchItem{В}{мають лише дві спільні точки.}
\matchItem{Г}{мають лише три спільні точки.}
\matchItem{Д}{мають чотири спільні точки.}
}{
\matchingGrid
}
% Відповідь: 1~--~Б, 2~--~А, 3~--~В
\ifshowsolutions\solution{\textbf{1}: прямі $y=2x-3$ і $y=x+1$ мають різні кутові коефіцієнти~--- одна спільна точка (Б). \textbf{2}: параболи $y=x^2$ і $y=x^2-2$ зсунуті по вертикалі~--- спільних точок немає (А). \textbf{3}: гіпербола $y=\dfrac{2}{x}$ і коло $x^2+y^2=4$ перетинаються у двох точках (В). Відповідь: \textbf{1~--~Б, 2~--~А, 3~--~В}.}\fi
\par\penalty-20
\end{samepage}

\begin{samepage}
\needspace{0.45\textheight}\typeTitle{Графік функції на рисунку: значення, нулі, точки екстремуму, зсуви}
\noindent\zadnum \begin{minipage}[t]{0.55\textwidth}
На рисунку зображено графік функції $y=f(x)$, визначеної на проміжку $[1; 9]$. Установіть відповідність між початком речення \mbox{(1--3)} та його закінченням \mbox{(А--Д)} так, щоб утворилося правильне твердження. \nmtyear{2024}
\end{minipage}
\hfill
\begin{minipage}[t]{0.4\textwidth}
\nbvspace{0pt}
\begin{flushright}
\begin{nmtfit}\begin{tikzpicture}[scale=0.5, thick]
    \draw[nmt grid, step=1] (-2,-2) grid (11,8);
    \draw[->] (-2.3,0) -- (11.3,0) node[right] {$x$};
    \draw[->] (0,-2) -- (0,8.3) node[above] {$y$};
    \node[below left] at (0,0) {$0$};
    \node[below] at (1,0) {$1$}; \node[below] at (9,0) {$9$};
    \node[left] at (0,1) {$1$};
    \draw[nmt curve] (1,5)
        .. controls (1.15,6.2) and (2,7) .. (3,7)
        .. controls (3.8,7) and (4.9,3) .. (5,0)
        .. controls (5.03,-0.9) and (5.5,-1) .. (6,-1)
        .. controls (6.5,-1) and (6.9,-0.5) .. (7,0)
        .. controls (7.1,0.5) and (8,1) .. (9,1);
    \node[nmt dot] at (1,5) {};
    \node[nmt dot] at (9,1) {};
    \node[right, fill=white, inner sep=1.5pt] at (4.0,7.3) {$y=f(x)$};
\end{tikzpicture}\end{nmtfit}
\end{flushright}
\end{minipage}

\nmtnobreak
\nopagebreak\nbvspace{0.3cm}

\nmtnobreak
\noindent
\matchingLayout{
\matchHead{Початок речення}
\matchItem{1}{Найбільше значення функції на проміжку $[1; 9]$}
\matchItem{2}{Найменше значення функції на проміжку $[1; 3]$}
\matchItem{3}{Найменше ціле значення $x$, за якого виконується нерівність $f(x) < 0$}
}{
\matchHead{Закінчення речення}
    \matchItem{А}{дорівнює $5$.}
\matchItem{Б}{дорівнює $6$.}
\matchItem{В}{дорівнює $7$.}
\matchItem{Г}{дорівнює $8$.}
\matchItem{Д}{дорівнює $9$.}
}{
\matchingGrid
}
\par\penalty-20\vspace{0.25cm}
\end{samepage}

\begin{samepage}
\noindent\zadnum \begin{minipage}[t]{0.55\textwidth}
На рисунку зображено графік функції $y=f(x)$, визначеної на проміжку $[-4; 4]$. Установіть відповідність між початком речення \mbox{(1--3)} та його закінченням \mbox{(А--Д)} так, щоб утворилося правильне твердження. \nmtyear{2024}
\end{minipage}
\hfill
\begin{minipage}[t]{0.4\textwidth}
\nbvspace{0pt}
\begin{flushright}
\begin{nmtfit}\begin{tikzpicture}[scale=0.5]
        \draw[nmt grid, step=1cm] (-4.5,-3.5) grid (4.5,3.5);
        \draw[->, thick] (-4.5,0) -- (4.5,0) node[below] {$x$};
        \draw[->, thick] (0,-3.5) -- (0,3.5) node[left] {$y$};

        \node[below left] at (0,0) {$0$};
        \node[below] at (1,0) {$1$};
        \node[left] at (0,1) {$1$};
        \node[below] at (4,0) {$4$};
        \node[below] at (-4,0) {$-4$};

        \draw[nmt curve] plot [smooth, tension=0.7] coordinates {(-4, -3) (-2, -1) (0, 0) (3, 3) (4, 1)};

        \path (-4,-3) node[nmt dot] {};
        \path (4,1) node[nmt dot] {};

        \node[left, fill=white, inner sep=1.5pt] at (-1, 1) {$y=f(x)$};
    \end{tikzpicture}\end{nmtfit}
\end{flushright}
\end{minipage}

\nmtnobreak
\nopagebreak\nbvspace{0.3cm}

\nmtnobreak
\noindent
\matchingLayout{
\matchHead{Початок речення}
\matchItem{1}{Найменше значення функції $y=f(x)$}
\matchItem{2}{Точка екстремуму функції $y=f(x)-5$}
\matchItem{3}{Нуль функції $y=f(x+2)$}
}{
\matchHead{Закінчення речення}
    \matchItem{А}{дорівнює $-3$.}
\matchItem{Б}{дорівнює $-2$.}
\matchItem{В}{дорівнює $0$.}
\matchItem{Г}{дорівнює $2$.}
\matchItem{Д}{дорівнює $3$.}
}{
\matchingGrid
}
% ... (попередній код)
\par\penalty-20\vspace{0.25cm}
\end{samepage}

\begin{samepage}
\noindent\zadnum На рисунку зображено графік функції $y=f(x)$, визначеної на проміжку $[-4; 4]$. Установіть відповідність між початком речення \mbox{(1--3)} та його закінченням \mbox{(А--Д)} так, щоб утворилося правильне твердження. \nmtyear{2024}

\nmtnobreak
\nopagebreak\nbvspace{0.3cm}

\nmtnobreak
% --- ВЕРХНІЙ БЛОК: Умови + Графік ---
\noindent
\begin{minipage}[t]{0.55\textwidth}
    \textit{Початок речення} \par\par\nbvspace{0.3cm}
    \textbf{1} \quad Найменше значення функції $y=f(x)$ \\[0.4cm]
    \textbf{2} \quad Точка екстремуму функції $y=f(x)-5$ \\[0.4cm]
    \textbf{3} \quad Нуль функції $y=f(x+2)$
\end{minipage}%
\hfill
\begin{minipage}[t]{0.40\textwidth}
    \nopagebreak\nbvspace{-0.3cm} % Підтягуємо графік трохи вгору
    \begin{flushright}
    \begin{nmtfit}\begin{tikzpicture}[scale=0.5]
        % Сітка та осі
        \draw[nmt grid, step=1cm] (-4.5,-3.5) grid (4.5,3.5);
        \draw[->, thick] (-4.5,0) -- (4.5,0) node[below] {$x$};
        \draw[->, thick] (0,-3.5) -- (0,3.5) node[left] {$y$};

        % Підписи осей
        \node[below left] at (0,0) {$0$};
        \node[below] at (1,0) {$1$};
        \node[left] at (0,1) {$1$};
        \node[below] at (4,0) {$4$};
        \node[below] at (-4,0) {$-4$};

        % Графік: старт (-4, -3), через (-2, -1), нуль (0,0), макс (3, 3), кінець (4, 1)
        \draw[nmt curve] plot [smooth, tension=0.7] coordinates {(-4, -3) (-2, -1) (0, 0) (3, 3) (4, 1)};

        % Точки на кінцях
        \path (-4,-3) node[nmt dot] {};
        \path (4,1) node[nmt dot] {};

        \node[left, fill=white, inner sep=1.5pt] at (-1, 1) {$y=f(x)$};
    \end{tikzpicture}\end{nmtfit}
    \end{flushright}
\end{minipage}

\nmtnobreak
\nopagebreak\nbvspace{0.2cm}

\nmtnobreak
% --- НИЖНІЙ БЛОК: Варіанти + Таблиця ---
\noindent
\begin{minipage}[t]{0.55\textwidth}
    \textit{Закінчення речення} \par\par\nbvspace{0.2cm}
    \begin{tabular}{ll}
    \textbf{А} & дорівнює $-3$. \\
    \textbf{Б} & дорівнює $-2$. \\
    \textbf{В} & дорівнює $0$. \\
    \textbf{Г} & дорівнює $2$. \\
    \textbf{Д} & дорівнює $3$. \\
    \end{tabular}
\end{minipage}%
\hfill
\begin{minipage}[t]{0.40\textwidth}
    \nopagebreak\nbvspace{0.5cm} % Вирівнювання таблички по висоті
    \begin{flushright}
    % Табличка
    \begingroup
    \setlength{\tabcolsep}{4pt}
    \renewcommand{\arraystretch}{1.2}
    \small
    \begin{tabular}{r|c|c|c|c|c|}
         \multicolumn{1}{c}{} & \multicolumn{1}{c}{\textbf{А}} & \multicolumn{1}{c}{\textbf{Б}} & \multicolumn{1}{c}{\textbf{В}} & \multicolumn{1}{c}{\textbf{Г}} & \multicolumn{1}{c}{\textbf{Д}} \\ \cline{2-6}
         \textbf{1} & & & & & \\ \cline{2-6}
         \textbf{2} & & & & & \\ \cline{2-6}
         \textbf{3} & & & & & \\ \cline{2-6}
    \end{tabular}
    \endgroup
    \end{flushright}
\end{minipage}
\par\penalty-20\vspace{0.25cm}
\end{samepage}

\begin{samepage}
\needspace{0.45\textheight}\typeTitle{Графік із кількох частин: проміжок --- формула}
\noindent\zadnum На рисунку зображено графік функції $y=f(x)$, визначеної на проміжку $[-4; 5]$. Узгодьте проміжок (1–3) з властивістю (А – Д), яку має функція на цьому проміжку. \nmtyear{2025}

\nmtnobreak
\nopagebreak\nbvspace{0.2cm}

\nmtnobreak
\noindent
\begin{minipage}[t]{0.55\textwidth}
    \nopagebreak\nbvspace{0pt}
\matchHead{Проміжок}
    \matchItem{1}{$[-4; -2]$}
\matchItem{2}{$[-2; 2]$}
\matchItem{3}{$[2; 5]$}

    \nopagebreak\nbvspace{0.3cm}
\matchHead{Властивість функції}
    \matchItem{А}{функція спадає}
\matchItem{Б}{функція має точку максимуму}
\matchItem{В}{графік функції перетинає пряму $y=-4{,}5$}
\matchItem{Г}{для кожної точки $(x_0; y_0)$, що належить графіку функції, добуток $x_0 \cdot y_0 < 0$}
\matchItem{Д}{функція набуває лише додатних значень}

    \nopagebreak\nbvspace{0.3cm}

    \matchingGrid
\end{minipage}%
\hfill
\begin{minipage}[t]{0.40\textwidth}
    \nopagebreak\nbvspace{0pt}
    \begin{flushright}
    \begin{nmtfit}\begin{tikzpicture}[scale=0.5]
        \draw[nmt grid, step=1cm] (-6.5,-7.5) grid (5.5,4.5);
        \draw[->, thick] (-6.5,0) -- (5.5,0) node[below] {$x$};
        \draw[->, thick] (0,-7.5) -- (0,4.5) node[left] {$y$};

        \node[below left] at (0,0) {$0$};
        \node[above] at (1,0) {$1$};
        \node[above left] at (0,1) {$1$};
        \node[below] at (5,0) {$5$};
        \node[below] at (-4,0) {$-4$};

        % Графік (хвиля)
        \draw[nmt curve] plot [smooth, tension=0.44] coordinates {(-4, -6) (-3, -5) (-2, -3)(-1, 0) (0, 1)(1, 0) (2, -3) (3, -4) (5, -1)};
        \path (-4,-6) node[nmt dot] {};
        \path (5, -1) node[nmt dot] {};
        \node[below] at (3, -3.5) {\normalsize $y=f(x)$};
    \end{tikzpicture}\end{nmtfit}
    \end{flushright}
\end{minipage}
\par\penalty-20\vspace{0.25cm}
\end{samepage}

\begin{samepage}
\noindent\zadnum На рисунку зображено графік функції $y=f(x)$, визначеної на проміжку $[-4; 5]$. Узгодьте проміжок (1–3) з властивістю (А – Д), яку має функція на цьому проміжку. \nmtyear{2025}

\nmtnobreak
\nopagebreak\nbvspace{0.2cm}

\nmtnobreak
\noindent
\begin{minipage}[t]{0.55\textwidth}
    \nopagebreak\nbvspace{0pt}
\matchHead{Проміжок}
    \matchItem{1}{$[-4; -2]$}
\matchItem{2}{$[-2; 2]$}
\matchItem{3}{$[2; 5]$}

    \nopagebreak\nbvspace{0.3cm}
\matchHead{Властивість функції}
    \matchItem{А}{функція зростає}
\matchItem{Б}{функція має точку мінімуму}
\matchItem{В}{функція спадає}
\matchItem{Г}{графік функції є фрагментом графіка функції $y=1-x^2$}
\matchItem{Д}{функція набуває лише додатних значень}

    \nopagebreak\nbvspace{0.3cm}

    \matchingGrid
\end{minipage}%
\hfill
\begin{minipage}[t]{0.40\textwidth}
    \nopagebreak\nbvspace{0pt}
    \begin{flushright}
    \begin{nmtfit}\begin{tikzpicture}[scale=0.5]
    \draw[nmt grid, step=1] (-5,-7) grid (6,3);
    \draw[->, thick] (-5,0) -- (6.4,0) node[right] {$x$};
    \draw[->, thick] (0,-7) -- (0,3.4) node[left] {$y$};
    \draw[thin] (-4,0.12) -- (-4,-0.12); \draw[thin] (1,0.12) -- (1,-0.12); \draw[thin] (5,0.12) -- (5,-0.12); \draw[thin] (0.12,1) -- (-0.12,1);
    \node[below left] at (0,0) {$0$};
    \node[below] at (-4,0) {$-4$};
    \node[below left] at (1,0) {$1$};
    \node[below left] at (5,0) {$5$};
    \node[anchor=south east, yshift=3pt] at (0,1) {$1$};
    \draw[nmt curve] (-4,-6) .. controls (-3.667,-5.833) and (-3.333,-5.444) .. (-3,-5) .. controls (-2.667,-4.556) and (-2.333,-3.8) .. (-2,-3) .. controls (-1.667,-2.2) and (-1.333,-0.5) .. (-1,0) .. controls (-0.667,0.5) and (-0.333,1) .. (0,1) .. controls (0.333,1) and (0.667,0.5) .. (1,0) .. controls (1.333,-0.5) and (1.667,-2.5) .. (2,-3) .. controls (2.333,-3.5) and (2.667,-4) .. (3,-4) .. controls (3.667,-4) and (4.333,-3.111) .. (5,-1);
    \node[nmt dot] at (-4,-6) {};
    \node[nmt dot] at (5,-1) {};
    \node[fill=white, inner sep=1.5pt] at (3,-5.3) {$y=f(x)$};
\end{tikzpicture}\end{nmtfit}
    \end{flushright}
\end{minipage}
\par\penalty-20\vspace{0.25cm}
\end{samepage}

\begin{samepage}
\needspace{0.45\textheight}\sectionTitle{Частина 2. З тригонометрією, показниковими функціями, логарифмами чи похідною}
\typeTitle{Функція --- її властивість:\\ парність, монотонність, період, область визначення}
\noindent\zadnum Установіть відповідність між функцією (1–3) та її властивістю (А–Д). \nmtyear{2023}\par
\nopagebreak\nbvspace{0.3cm}

\nmtnobreak
\noindent
\matchingLayout{
\matchHead{Функція}
    \matchItem{1}{$y = (x + 2)^2$}
\matchItem{2}{$y = 2\sqrt{x}$}
\matchItem{3}{$y = 2^x$}
}{
\matchHead{Властивість}
    \matchItem{А}{є зростаючою на $(-\infty; +\infty)$}
\matchItem{Б}{графік має одну спільну точку із графіком $y = x - 2$}
\matchItem{В}{графік має дві спільні точки із графіком $y = x - 2$}
\matchItem{Г}{є спадною на проміжку $(-\infty; -2]$}
\matchItem{Д}{є спадною на проміжку $(-\infty; 2]$}
}{
\matchingGrid
}
\par\penalty-20\vspace{0.25cm}
\end{samepage}

\begin{samepage}
\noindent\zadnum До кожного початку речення \mbox{(1--3)} доберіть його закінчення \mbox{(А--Д)} так, щоб утворилося правильне твердження. \nmtyear{2023}

\nmtnobreak
\nopagebreak\nbvspace{0.3cm}

\nmtnobreak
\noindent
\matchingLayout{
\matchHead{Початок речення}
\matchItem{1}{Функція $y = \log_2 x$}
\matchItem{2}{Функція $y = x^2 - 4x - 4$}
\matchItem{3}{Функція $y = \dfrac{1}{x}$}

\nmtnobreak
    % Зменшена табличка
}{
\matchHead{Закінчення речення}
    \matchItem{А}{не визначена при $x = -2$.}
\matchItem{Б}{набуває від’ємного значення при $x = 2$.}
\matchItem{В}{є непарною.}
\matchItem{Г}{спадає на проміжку $(-\infty; 4]$.}
\matchItem{Д}{зростає на проміжку $(-\infty; +\infty)$.}
}{
\matchingGrid
}
\par\penalty-20\vspace{0.25cm}
\end{samepage}

\begin{samepage}
\noindent\zadnum Установіть відповідність між функцією \mbox{(1--3)} та її властивістю \mbox{(А--Д)}. \nmtyear{2023}

\nmtnobreak
\nopagebreak\nbvspace{0.3cm}

\nmtnobreak
\noindent
\matchingLayout{
\matchHead{Функція}
\matchItem{1}{$y = \sqrt{x+1}$}
\matchItem{2}{$y = 4-x^2$}
\matchItem{3}{$y = 3^{-x}$}

\nmtnobreak
    % Зменшена табличка
}{
\matchHead{Властивість}
    \matchItem{А}{має точку локального максимуму}
\matchItem{Б}{має точку локального мінімуму}
\matchItem{В}{є непарною}
\matchItem{Г}{зростає на всій області визначення}
\matchItem{Д}{набуває лише додатних значень}
}{
\matchingGrid
}
\par\penalty-20\vspace{0.25cm}
\end{samepage}

\begin{samepage}
\noindent\zadnum До кожного початку речення \mbox{(1--3)} доберіть його закінчення \mbox{(А--Д)} так, щоб утворилося правильне твердження. \nmtyear{2023}

\nmtnobreak
\nopagebreak\nbvspace{0.3cm}

\nmtnobreak
\noindent
\matchingLayout{
\matchHead{Початок речення}
\matchItem{1}{Функція $y = \log_{0{,}5} x$}
\matchItem{2}{Функція $y = \sin x$}
\matchItem{3}{Функція $y = \dfrac{1}{2x-2}$}

\nmtnobreak
    % Зменшена табличка
}{
\matchHead{Закінчення речення}
    \matchItem{А}{не визначена при $x=1$.}
\matchItem{Б}{набуває від’ємного значення при $x=2$.}
\matchItem{В}{є непарною.}
\matchItem{Г}{має лише одну точку локального екстремуму.}
\matchItem{Д}{зростає на проміжку $(0; +\infty)$.}
}{
\matchingGrid
}
\par\penalty-20\vspace{0.25cm}
\end{samepage}

\begin{samepage}
\noindent\zadnum Установіть відповідність між функцією \mbox{(1--3)} та її властивістю \mbox{(А--Д)}. \nmtyear{2023}

\nmtnobreak
\nopagebreak\nbvspace{0.3cm}

\nmtnobreak
\noindent
\matchingLayout{
\matchHead{Функція}
\matchItem{1}{$y = 7x + 4$}
\matchItem{2}{$y = -\dfrac{7}{x}$}
\matchItem{3}{$y = \log_{0{,}5} (x-4)$}

\nmtnobreak
    % Зменшена табличка
}{
\matchHead{Властивість}
    \matchItem{А}{є спадною на всій області визначення}
\matchItem{Б}{графік функції перетинає вісь $y$ в точці з ординатою $4$}
\matchItem{В}{є непарною}
\matchItem{Г}{є парною}
\matchItem{Д}{областю визначення є проміжок $(0; +\infty)$}
}{
\matchingGrid
}
\par\penalty-20\vspace{0.25cm}
\end{samepage}

\begin{samepage}
\noindent\zadnum Установіть відповідність між функцією \mbox{(1--3)} та її властивістю \mbox{(А--Д)}. \nmtyear{2023}\par
\nopagebreak\nbvspace{0.3cm}

\nmtnobreak
\matchingLayout{
\textit{Функція}

\nmtnobreak
\matchItem{1}{$y = x^2 + 4$}
\matchItem{2}{$y = \log_4 x$}
\matchItem{3}{$y = 4 - 2x$}
}{
\matchHead{Властивість}

\nmtnobreak
\matchItem{А}{має лише одну точку локального екстремуму}
\matchItem{Б}{є зростаючою на всій області визначення}
\matchItem{В}{має від'ємний нуль}
\matchItem{Г}{є спадною на всій області визначення}
\matchItem{Д}{непарна}
}{
\matchingGrid
}
\par\penalty-20\vspace{0.25cm}
\end{samepage}

\begin{samepage}
\noindent\zadnum Установіть відповідність між функцією \mbox{(1--3)} та її властивістю \mbox{(А--Д)}. \nmtyear{2024}

\nmtnobreak
\nopagebreak\nbvspace{0.3cm}

\nmtnobreak
\noindent
\matchingLayout{
\matchHead{Функція}
\matchItem{1}{$y = 4 - x^2$}
\matchItem{2}{$y = -x^3$}
\matchItem{3}{$y = 4^x$}

\nmtnobreak
    % Табличка
}{
\matchHead{Властивість}
    \matchItem{А}{є зростаючою на всій області визначення}
\matchItem{Б}{набуває від’ємного значення при $x = -2$}
\matchItem{В}{є парною}
\matchItem{Г}{має три спільні точки з графіком функції $y = -x$}
\matchItem{Д}{графік функції розташований лише в I та IV координатній чверті}
}{
\matchingGrid
}
\par\penalty-20\vspace{0.25cm}
\end{samepage}

\begin{samepage}
\noindent\zadnum Узгодьте функцію (1–3) із її властивістю (А – Д). \nmtyear{2025}

\nmtnobreak
\nopagebreak\nbvspace{0.3cm}

\nmtnobreak
\noindent
\matchingLayout{
\matchHead{Функція}
    \matchItem{1}{$y=x^2$}
\matchItem{2}{$y=x^3+1$}
\matchItem{3}{$y=\log_2 x$}
}{
\matchHead{Властивість функції}
    \matchItem{А}{парна}
\matchItem{Б}{непарна}
\matchItem{В}{область визначення функції є проміжок $(0; +\infty)$}
\matchItem{Г}{похідна функції є додатною на проміжку $(-\infty; 0)$}
\matchItem{Д}{є спадною на всій області визначення}
}{
\matchingGrid
}
\par\penalty-20\vspace{0.25cm}
\end{samepage}

\begin{samepage}
\noindent\zadnum Доберіть до функції \mbox{(1--3)} її властивість \mbox{(А--Д)}. \nmtyear{2025}

\nmtnobreak
\nopagebreak\nbvspace{0.3cm}

\nmtnobreak
\noindent
\matchingLayout{
\matchHead{Функція}
    \matchItem{1}{$y=2x-5$}
\matchItem{2}{$y=\log_3(x-5)$}
\matchItem{3}{$y=x^2-5$}
}{
\matchHead{Властивість функції}
    \matchItem{А}{графік функції має лише 2 точки перетину з осями координат}
\matchItem{Б}{областю значень функції є проміжок $[-5; +\infty)$}
\matchItem{В}{графік функції знаходиться лише в I та II координатних чвертях}
\matchItem{Г}{областю визначення функції є проміжок $(5; +\infty)$}
\matchItem{Д}{функція спадає на всій області визначення}
}{
\matchingGrid
}
\par\penalty-20\vspace{0.25cm}
\end{samepage}

\begin{samepage}
\noindent\zadnum Увідповідніть функцію (1–3) та її властивість (А – Д). \nmtyear{2025}

\nmtnobreak
\nopagebreak\nbvspace{0.3cm}

\nmtnobreak
\noindent
\matchingLayout{
\matchHead{Функція}
    \matchItem{1}{$f(x)=x^2+2$}
\matchItem{2}{$f(x)=\sin x$}
\matchItem{3}{$f(x)=x^3$}
}{
\matchHead{Властивість функції}
    \matchItem{А}{зростає на всій області визначення}
\matchItem{Б}{спадає на всій області визначення}
\matchItem{В}{парна}
\matchItem{Г}{має більше ніж два нулі}
\matchItem{Д}{графік функції симетричний відносно осі $x$}
}{
\matchingGrid
}
\par\penalty-20\vspace{0.25cm}
\end{samepage}

\begin{samepage}
\noindent\zadnum Узгодьте функцію (1–3) із її властивістю (А – Д). \nmtyear{2025}

\nmtnobreak
\nopagebreak\nbvspace{0.3cm}

\nmtnobreak
\noindent
\matchingLayout{
\matchHead{Функція}
    % Тут формули короткі, тому p{5cm} вистачить
    \matchItem{1}{$y=-x^3$}
\matchItem{2}{$y=\sqrt{x}$}
\matchItem{3}{$y=\cos x$}
}{
\matchHead{Властивість функції}
    % Тут текст довгий, тому використовуємо p{6.2cm} для перенесення слів
    \matchItem{А}{є парною}
\matchItem{Б}{має лише одну спільну точку з колом $x^2+y^2=9$}
\matchItem{В}{немає спільних точок з прямою $x=0$}
\matchItem{Г}{графік функції розташований лише в другій координатній чверті}
\matchItem{Д}{набуває всіх значень з проміжку $(-\infty; +\infty)$}
}{
\matchingGrid
}
\par\penalty-20\vspace{0.25cm}
\end{samepage}

\begin{samepage}
\noindent\zadnum Узгодьте функцію (1–3) із її властивістю (А – Д). \nmtyear{2025}

\nmtnobreak
\nopagebreak\nbvspace{0.3cm}

\nmtnobreak
\noindent
\matchingLayout{
\matchHead{Функція}
    \matchItem{1}{$y=7x+4$}
\matchItem{2}{$y=-\dfrac{7}{x}$}
\matchItem{3}{$y=\log_{0{,}1}(x-4)$}
}{
\matchHead{Властивість функції}
    \matchItem{А}{парна}
\matchItem{Б}{непарна}
\matchItem{В}{область визначення функції є проміжок $(0; +\infty)$}
\matchItem{Г}{графік функції перетинає вісь $y$ в точці з ординатою $4$}
\matchItem{Д}{спадає на всій області визначення}
}{
\matchingGrid
}
\par\penalty-20\vspace{0.25cm}
\end{samepage}

\begin{samepage}
\noindent\zadnum Узгодьте функцію (1–3) із її властивістю (А – Д). \nmtyear{2025}

\nmtnobreak
\nopagebreak\nbvspace{0.3cm}

\nmtnobreak
\noindent
\matchingLayout{
\matchHead{Функція}
    \matchItem{1}{$y=\lg x$}
\matchItem{2}{$y=\dfrac{1}{x}$}
\matchItem{3}{$y=(x+1)^2$}
}{
\matchHead{Властивість функції}
    \matchItem{А}{зростає на всій області визначення}
\matchItem{Б}{парна}
\matchItem{В}{має точку максимуму}
\matchItem{Г}{не визначена лише в одній точці}
\matchItem{Д}{має точку мінімуму}
}{
\matchingGrid
}
\par\penalty-20\vspace{0.25cm}
\end{samepage}

\begin{samepage}
% НМТ 2026, сесія 12.06, завдання №17 --- також у темі: Показникова функція і показникові вирази
\zadtask{Доберіть до функції \mbox{(1--3)} її властивість \mbox{(А--Д)}. \nmtyear{2026}\par
\nopagebreak\nbvspace{0.3cm}
\noindent
\matchingLayout{
\matchHead{Функція}
\matchItem{1}{$y = \sqrt{x - 3}$}
\matchItem{2}{$y = x^2 + 3$}
\matchItem{3}{$y = 3^x$}
}{
\matchHead{Властивість функції}
\matchItem{А}{областю значень є проміжок $[3; +\infty)$}
\matchItem{Б}{областю визначення є проміжок $[3; +\infty)$}
\matchItem{В}{областю визначення є проміжок $[-3; +\infty)$}
\matchItem{Г}{зростає на множині всіх дійсних чисел}
\matchItem{Д}{є непарною}
}{
\matchingGrid
}
}
\par\penalty-20
\end{samepage}

\begin{samepage}
% НМТ 2026, сесія 5.06, завдання №17 --- також у темі: Логарифмічна функція
\zadnum Доберіть до функції \mbox{(1--3)} її властивість \mbox{(А--Д)}. \nmtyear{2026}

\nmtnobreak
\nopagebreak\nbvspace{0.3cm}
\noindent
\matchingLayout{
\matchHead{Функція}
\matchItem{1}{$y = x^4 - x^2$}
\matchItem{2}{$y = \log_{0{,}2} x$}
\matchItem{3}{$y = -\dfrac{1}{x}$}
}{
\matchHead{Властивість функції}
\matchItem{А}{область визначення збігається з областю значень}
\matchItem{Б}{є спадною}
\matchItem{В}{має лише три нулі}
\matchItem{Г}{має лише два нулі}
\matchItem{Д}{графік розташований у I та III чвертях}
}{
\matchingGrid
}
% Відповідь: 1~--~В, 2~--~Б, 3~--~А
\ifshowsolutions\solution{$1$: $y=x^4-x^2=x^2(x^2-1)$ має нулі $x=0;\,\pm1$, тобто три нулі (В); $2$: $y=\log_{0{,}2}x$ при основі $0<0{,}2<1$ спадна (Б); $3$: $y=-\dfrac1x$ має область визначення $x\ne0$, що збігається з областю значень (А). Відповідь: \textbf{1~--~В, 2~--~Б, 3~--~А}.}\fi
\par\penalty-20
\end{samepage}

\begin{samepage}
% НМТ 2026, сесія 8.06, завдання №17 --- також у темі: Показникова функція і показникові вирази, Логарифмічна функція
\zadtask{Доберіть до функції \mbox{(1--3)} її властивість \mbox{(А--Д)}. \nmtyear{2026}\par
\nopagebreak\nbvspace{0.3cm}
\noindent
\matchingLayout{
\matchHead{Функція}
\matchItem{1}{$y = \lg x$}
\matchItem{2}{$y = 0{,}2^x - 1$}
\matchItem{3}{$y = x^2 - 4x + 5$}
}{
\matchHead{Властивість}
\matchItem{А}{є зростаючою}
\matchItem{Б}{є спадною}
\matchItem{В}{не має нулів}
\matchItem{Г}{має лише два нулі}
\matchItem{Д}{є парною}
}{
\matchingGrid
}
}
\par\penalty-20
\end{samepage}

\begin{samepage}
% НМТ 2026, сесія 9.06, завдання №17
\zadtask{Доберіть до функції \mbox{(1--3)} її властивість \mbox{(А--Д)}. \nmtyear{2026}\par
\nopagebreak\nbvspace{0.3cm}
\noindent
\matchingLayout{
\matchHead{Функція}
\matchItem{1}{$y = \sin x$}
\matchItem{2}{$y = x^2 + 2$}
\matchItem{3}{$y = 3x + 2$}
}{
\matchHead{Властивість функції}
\matchItem{А}{парна}
\matchItem{Б}{графік функції симетричний відносно осі $x$}
\matchItem{В}{спадає на області визначення}
\matchItem{Г}{зростає на області визначення}
\matchItem{Д}{графік функції проходить через точку $(0; 0)$}
}{
\matchingGrid
}
}
\par\penalty-20
\end{samepage}

\begin{samepage}
\needspace{0.45\textheight}\typeTitle{Графік функції --- його властивість: перетин з осями, симетрія}
\noindent\zadnum Установіть відповідність між функцією (1–3) та властивістю (А – Д) її графіка. \nmtyear{2023}

\nmtnobreak
\nopagebreak\nbvspace{0.3cm}

\nmtnobreak
\noindent
\matchingLayout{
\matchHead{Функція}
    \matchItem{1}{$y=\cos x$}
\matchItem{2}{$y=-\dfrac{2}{x}$}
\matchItem{3}{$y=2^x$}
}{
\matchHead{Властивість графіка}
    \matchItem{А}{проходить через точку $(1; 0)$}
\matchItem{Б}{не перетинає вісь $y$}
\matchItem{В}{симетричний відносно осі $y$}
\matchItem{Г}{розміщений лише в I та II чвертях}
\matchItem{Д}{розміщений лише в I та IV чвертях}
}{
\matchingGrid
}
\par\penalty-20\vspace{0.25cm}
\end{samepage}

\begin{samepage}
\noindent\zadnum Установіть відповідність між функцією (1–3) та властивістю (А – Д) її графіка. \nmtyear{2023}

\nmtnobreak
\nopagebreak\nbvspace{0.3cm}

\nmtnobreak
\noindent
\matchingLayout{
\matchHead{Функція}
    \matchItem{1}{$y=\dfrac{4}{x}$}
\matchItem{2}{$y=\sin x$}
\matchItem{3}{$y=\log_2 x$}
}{
\matchHead{Властивість графіка функції}
    \matchItem{А}{не перетинає вісь $x$}
\matchItem{Б}{перетинає вісь $x$ у точці з абсцисою $1$}
\matchItem{В}{двічі перетинає графік функції $y=(x-1)^2$}
\matchItem{Г}{симетричний відносно осі $y$}
\matchItem{Д}{розміщений лише в першій і другій координатних чвертях}
}{
\matchingGrid
}
\par\penalty-20\vspace{0.25cm}
\end{samepage}

\begin{samepage}
\noindent\zadnum Установіть відповідність між функцією (1–3) та її властивістю (А – Д). \nmtyear{2023}

\nmtnobreak
\nopagebreak\nbvspace{0.3cm}

\nmtnobreak
\noindent
\matchingLayout{
\matchHead{Функція}
    \matchItem{1}{$y=x^2-1$}
\matchItem{2}{$y=3^x$}
\matchItem{3}{$y=x^3$}
}{
\matchHead{Властивість}
    \matchItem{А}{проходить через початок координат}
\matchItem{Б}{не містить точок з від'ємною ординатою}
\matchItem{В}{не має спільних точок з віссю $y$}
\matchItem{Г}{двічі перетинає графік прямої $y=3$}
\matchItem{Д}{графік знаходиться лише в I чверті}
}{
\matchingGrid
}
\par\penalty-20\vspace{0.25cm}
\end{samepage}

\begin{samepage}
\noindent\zadnum Установіть відповідність між функцією \mbox{(1--3)} та властивістю її графіка \mbox{(А--Д)}. \nmtyear{2024}

\nmtnobreak
\nopagebreak\nbvspace{0.3cm}

\nmtnobreak
\noindent
\matchingLayout{
\matchHead{Функція}
    \matchItem{1}{$y=x+2$}
\matchItem{2}{$y=x$}
\matchItem{3}{$y=4$}
}{
\matchHead{Властивість графіка функції}
    \matchItem{А}{спадає}
\matchItem{Б}{утворює з осями координат рівнобедрений трикутник}
\matchItem{В}{немає спільних із графіком функції $y=\log_{0{,}5} x$}
\matchItem{Г}{перетинає графік рівняння $x^2 + y^2 = 1$}
\matchItem{Д}{не перетинає вісь абсцис}
}{
\matchingGrid
}
\par\penalty-20\vspace{0.25cm}
\end{samepage}

\begin{samepage}
\noindent\zadnum Установіть відповідність між функцією \mbox{(1--3)} та її властивістю її графіка \mbox{(А--Д)}. \nmtyear{2024}

\nmtnobreak
\nopagebreak\nbvspace{0.3cm}

\nmtnobreak
\noindent
\matchingLayout{
\matchHead{Функція}
    \matchItem{1}{$y=2x^3$}
\matchItem{2}{$y=\dfrac{2}{x}-1$}
\matchItem{3}{$y=\cos x$}
}{
\matchHead{Властивість графіка функції}
    \matchItem{А}{двічі перетинає пряму $x=1$}
\matchItem{Б}{симетричний відносно початку координат}
\matchItem{В}{симетричний відносно осі абсцис}
\matchItem{Г}{симетричний відносно осі ординат}
\matchItem{Д}{не перетинає вісь ординат}
}{
\matchingGrid
}
\par\penalty-20\vspace{0.25cm}
\end{samepage}

\begin{samepage}
\noindent\zadnum Установіть відповідність між функцією \mbox{(1--3)} та властивістю її графіка \mbox{(А--Д)}. \nmtyear{2024}

\nmtnobreak
\nopagebreak\nbvspace{0.3cm}

\nmtnobreak
\noindent
\matchingLayout{
\matchHead{Функція}
    \matchItem{1}{$y=\sqrt{x+4}$}
\matchItem{2}{$y=\dfrac{4}{x}$}
\matchItem{3}{$y=4^x$}
}{
\matchHead{Властивість графіка функції}
    \matchItem{А}{розташований лише в першій чверті}
\matchItem{Б}{має спільну точку з віссю $x$}
\matchItem{В}{проходить через точку $(0; 1)$}
\matchItem{Г}{не перетинає вісь $y$}
\matchItem{Д}{симетричний відносно осі $y$}
}{
\matchingGrid
}
\par\penalty-20\vspace{0.25cm}
\end{samepage}

\begin{samepage}
\noindent\zadnum До кожного початку речення \mbox{(1--3)} доберіть його закінчення \mbox{(А--Д)} так, щоб утворилося правильне твердження. \nmtyear{2024}

\nmtnobreak
\nopagebreak\nbvspace{0.3cm}

\nmtnobreak
\noindent
\matchingLayout{
\matchHead{Початок речення}
    \matchItem{1}{Графік функції $y=2x$}
\matchItem{2}{Графік функції $y=\log_2 x$}
\matchItem{3}{Графік функції $y=2^x$}
}{
\matchHead{Закінчення речення}
    \matchItem{А}{симетричний відносно осі абсцис.}
\matchItem{Б}{симетричний відносно осі ординат.}
\matchItem{В}{симетричний відносно початку координат.}
\matchItem{Г}{не перетинає вісь абсцис.}
\matchItem{Д}{не перетинає вісь ординат.}
}{
\matchingGrid
}
\par\penalty-20\vspace{0.25cm}
\end{samepage}

\begin{samepage}
\noindent\zadnum Установіть відповідність між функцією \mbox{(1--3)} та властивістю її графіка \mbox{(А--Д)}. \nmtyear{2024}

\nmtnobreak
\nopagebreak\nbvspace{0.3cm}

\nmtnobreak
\noindent
\matchingLayout{
\matchHead{Функція}
    \matchItem{1}{$y=2x+6$}
\matchItem{2}{$y=-2$}
\matchItem{3}{$y=-x$}
}{
\matchHead{Властивість графіка функції}
    \matchItem{А}{паралельний осі $y$}
\matchItem{Б}{є бісектрисою другої та четвертої координатних чвертей}
\matchItem{В}{дотикається до графіка рівняння $x^2+y^2=4$}
\matchItem{Г}{паралельний до графіка функції $y=2x$}
\matchItem{Д}{не перетинає графік функції $y=\tg x$}
}{
\matchingGrid
}
\par\penalty-20\vspace{0.25cm}
\end{samepage}

\begin{samepage}
\noindent\zadnum Доберіть до функції \mbox{(1--3)} властивість її графіка \mbox{(А--Д)}. \nmtyear{2025}

\nmtnobreak
\nopagebreak\nbvspace{0.3cm}

\nmtnobreak
\noindent
\matchingLayout{
\matchHead{Функція}
    \matchItem{1}{$y=\dfrac{1}{x}-1$}
\matchItem{2}{$y=x^2-4x+5$}
\matchItem{3}{$y=2\cos x$}
}{
\matchHead{Властивість графіка функції}
    \matchItem{А}{перетинає лише вісь $y$}
\matchItem{Б}{перетинає лише вісь $x$}
\matchItem{В}{не перетинає жодну з осей координат}
\matchItem{Г}{симетричний відносно осі $x$}
\matchItem{Д}{симетричний відносно осі $y$}
}{
\matchingGrid
}
\par\penalty-20\vspace{0.25cm}
\end{samepage}

\begin{samepage}
% НМТ 2026, сесія 1.06, завдання №17 --- основна тема: Графіки елементарних функцій, їх побудова та перетворення
\zadtask{Увідповідніть початок речення \mbox{(1--3)} із його закінченням \mbox{(А--Д)} так, щоб утворилося правильне твердження. \nmtyear{2026}

\nmtnobreak
\nopagebreak\nbvspace{0.2cm}
\noindent
\matchingLayout{
\matchHead{Початок речення}
\matchItem{1}{Графік функції $y = \dfrac{2}{x} + 3$}
\matchItem{2}{Графік функції $y = 4^x$}
\matchItem{3}{Графік функції $y = \mathrm{tg}\,x$}
}{
\matchHead{Закінчення речення}
\matchItem{А}{симетричний відносно початку координат.}
\matchItem{Б}{симетричний відносно осі $y$.}
\matchItem{В}{перетинає лише вісь $x$.}
\matchItem{Г}{перетинає лише вісь $y$.}
\matchItem{Д}{перетинає графік функції $y = x$ лише в одній точці.}
}{
\matchingGrid
}
}
% Відповідь: 1~--~В, 2~--~Г, 3~--~А
\ifshowsolutions\solution{\textbf{1)} $y=\dfrac{2}{x}+3$ ($x\ne0$): осі $y$ не перетинає, а вісь $x$ перетинає при $x=-\dfrac{2}{3}$ --- лише вісь $x$ (В); \textbf{2)} $y=4^x>0$: перетинає лише вісь $y$ у точці $(0;1)$ (Г); \textbf{3)} $y=\mathrm{tg}\,x$ --- непарна, симетрична відносно початку координат (А). Відповідь: \textbf{1\,--\,В, 2\,--\,Г, 3\,--\,А}.}\fi
\par\penalty-20
\end{samepage}

\begin{samepage}
% НМТ 2026, сесія 16.06, завдання №17 --- основна тема: Графіки елементарних функцій, їх побудова та перетворення
\zadtask{Доберіть до функції \mbox{(1--3)} властивість її графіка \mbox{(А--Д)}. \nmtyear{2026}\par
\nopagebreak\nbvspace{0.3cm}
\noindent
\matchingLayout{
\matchHead{Функція}
\matchItem{1}{$y = \sqrt{x-2} + 1$}
\matchItem{2}{$y = -5^x$}
\matchItem{3}{$y = x^3$}
}{
\matchHead{Властивість графіка функції}
\matchItem{А}{симетричний відносно початку координат}
\matchItem{Б}{симетричний відносно осі $y$}
\matchItem{В}{розміщений лише в ІІІ та IV координатних чвертях}
\matchItem{Г}{розміщений лише в І координатній чверті}
\matchItem{Д}{проходить через точку $(0; 1)$}
}{
\matchingGrid
}
}
% Відповідь: 1--Г; 2--В; 3--А
\par\penalty-20
\end{samepage}

\begin{samepage}
% НМТ 2026, сесія 17.06, завдання №17 --- також у темі: Логарифмічна функція
\zadtask{Доберіть до функції \mbox{(1--3)} властивість її графіка \mbox{(А--Д)}. \nmtyear{2026}\par
\nopagebreak\nbvspace{0.3cm}
\noindent
\matchingLayout{
\matchHead{Функція}
\matchItem{1}{$y = 2 - x$}
\matchItem{2}{$y = \dfrac{1}{x - 2}$}
\matchItem{3}{$y = \log_{0{,}5} x$}
}{
\matchHead{Властивість графіка функції}
\matchItem{А}{перетинає вісь $y$ з додатною ординатою, а вісь $x$ не перетинає}
\matchItem{Б}{не перетинає ані вісь $x$, ані вісь $y$}
\matchItem{В}{розміщений лише в I, II та IV координатних чвертях}
\matchItem{Г}{перетинає вісь $y$ з від'ємною ординатою, а вісь $x$ не перетинає}
\matchItem{Д}{перетинає вісь $x$, а вісь $y$ не перетинає}
}{
\matchingGrid
}
}
% Відповідь: 1--В, 2--Г, 3--Д
\par\penalty-20
\end{samepage}

\begin{samepage}
% НМТ 2026, сесія 15.06, завдання №17 --- також у темі: Показникова функція і показникові вирази
\zadtask{До кожного початку речення \mbox{(1--3)} доберіть його закінчення \mbox{(А--Д)} так, щоб утворилося правильне твердження. \nmtyear{2026}\par
\nopagebreak\nbvspace{0.3cm}
\noindent
\matchingLayout{
\matchHead{Початок речення}
\matchItem{1}{Графіки функцій $y = 2^x$ та $y = 0{,}5^x$ перетинаються в точці}
\matchItem{2}{Пряма $y = x - 1$ перетинає вісь $x$ у точці}
\matchItem{3}{Графік функції $y = (x + 2)^2$ дотикається осі $x$ у точці}
}{
\matchHead{Закінчення речення}
\matchItem{А}{$(1; 0)$.}
\matchItem{Б}{$(0; 1)$.}
\matchItem{В}{$(2; 0)$.}
\matchItem{Г}{$(-2; 0)$.}
\matchItem{Д}{$(-1; 0)$.}
}{
\matchingGrid
}
}
% Відповідь: 1--Б, 2--А, 3--Г
\par\penalty-20
\end{samepage}

\begin{samepage}
\needspace{0.45\textheight}\typeTitle{Твердження --- функція, для якої воно правильне}
\noindent\zadnum Установіть відповідність між твердженням \mbox{(1--3)} та функцією \mbox{(А--Д)}, для якої це твердження є правильним. \nmtyear{2023}

\nmtnobreak
\nopagebreak\nbvspace{0.3cm}

\nmtnobreak
\noindent
\matchingLayout{
% Трохи ширше для довгих тверджень
\matchHead{Твердження}
    \matchItem{1}{областю визначення функції є проміжок $[-1; +\infty)$}
\matchItem{2}{графік функції проходить через точку $(0; 1)$}
\matchItem{3}{функція має точку локального екстремуму на проміжку $[1; 3]$}

\nmtnobreak
    % Зменшена табличка
}{
\matchHead{Функція}
    \matchItem{А}{$y = -\cos x$}
\matchItem{Б}{$y = 4^x$}
\matchItem{В}{$y = 2\sqrt{x+1}$}
\matchItem{Г}{$y = x^2 - 4x + 3$}
\matchItem{Д}{$y = x$}
}{
\matchingGrid
}
\par\penalty-20\vspace{0.25cm}
\end{samepage}

\begin{samepage}
\noindent\zadnum Установіть відповідність між твердженням \mbox{(1--3)} та функцією \mbox{(А--Д)}, для якої це твердження є правильним. \nmtyear{2023}

\nmtnobreak
\nopagebreak\nbvspace{0.3cm}

\nmtnobreak
\noindent
\matchingLayout{
\matchHead{Твердження}
    \matchItem{1}{областю значень функції є проміжок $[0; +\infty)$}
\matchItem{2}{графік функції симетричний відносно осі $y$}
\matchItem{3}{найменше значення на відрізку $[1; 4]$ функція набуває в точці $x = 4$}

\nmtnobreak
    % Зменшена табличка
}{
\matchHead{Функція}
    \matchItem{А}{$y = x^2 + 4$}
\matchItem{Б}{$y = x$}
\matchItem{В}{$y = \sqrt{x}$}
\matchItem{Г}{$y = \log_{0{,}5} x$}
\matchItem{Д}{$y = -\dfrac{1}{x}$}
}{
\matchingGrid
}
\par\penalty-20\vspace{0.25cm}
\end{samepage}

\begin{samepage}
\noindent\zadnum Установіть відповідність між твердженням \mbox{(1--3)} та функцією \mbox{(А--Д)}, для якої це твердження є правильним. \nmtyear{2024}

\nmtnobreak
\nopagebreak\nbvspace{0.3cm}

\nmtnobreak
\noindent
\matchingLayout{
\matchHead{Твердження}
    \matchItem{1}{є непарною}
\matchItem{2}{зростає на відрізку $[1; 4]$}
\matchItem{3}{найбільше значення функції на відрізку $[1; 4]$ є від'ємним числом}
}{
\matchHead{Функція}
    \matchItem{А}{$y = \dfrac{1}{x+1}$}
\matchItem{Б}{$y = \sin x$}
\matchItem{В}{$y = x^2 - 1$}
\matchItem{Г}{$y = 0{,}5^x$}
\matchItem{Д}{$y = -\sqrt{x}$}
}{
\matchingGrid
}
\par\penalty-20\vspace{0.25cm}
\end{samepage}

\begin{samepage}
\noindent\zadnum Установіть відповідність між твердженням \mbox{(1--3)} та функцією \mbox{(А--Д)}, для якої це твердження є правильним. \nmtyear{2024}

\nmtnobreak
\nopagebreak\nbvspace{0.3cm}

\nmtnobreak
\noindent
\matchingLayout{
\matchHead{Твердження}
    \matchItem{1}{функція має 2 нулі}
\matchItem{2}{на відрізку $[-1; 3]$ функція набуває від'ємних значень}
\matchItem{3}{найменше значення функції на відрізку $[-1; 3]$ дорівнює $0{,}5$}

\nmtnobreak
    % Табличка
}{
\matchHead{Функція}
    \matchItem{А}{$y = x^2 - 4$}
\matchItem{Б}{$y = \dfrac{1}{x-4}$}
\matchItem{В}{$y = 2^x$}
\matchItem{Г}{$y = 0{,}5^x$}
\matchItem{Д}{$y = \sqrt{x+1}$}
}{
\matchingGrid
}
\par\penalty-20\vspace{0.25cm}
\end{samepage}

\begin{samepage}
\needspace{0.45\textheight}\typeTitle{Координатні чверті, у яких лежить графік}
% перенесено з теми 26: завдання про графік або властивості функції
\noindent\zadnum Доберіть до функції \mbox{(1--3)} координатні чверті \mbox{(А--Д)}, у яких лежить графік цієї функції (координатні чверті показано на рисунку). \nmtyear{2025}

\nmtnobreak
\nopagebreak\nbvspace{0.3cm}

\nmtnobreak
\noindent
\matchingLayout{
\matchHead{Функція}
    \matchItem{1}{$y=-3-x^2$}
\matchItem{2}{$y=\dfrac{3}{x}$}
\matchItem{3}{\trigStyle{$y=\operatorname{tg} x$}}
}{
\matchHead{Координатні чверті}
    \matchItem{А}{лише I та II}
\matchItem{Б}{лише I та III}
\matchItem{В}{лише III та IV}
\matchItem{Г}{лише I та IV}
\matchItem{Д}{I, II, III та IV}
}{
\matchingGrid
}
\hfill
\begin{minipage}[t]{0.20\textwidth}
    \nopagebreak\nbvspace{0.5cm}
    \begin{nmtfit}\begin{tikzpicture}
    \draw[->] (-2.0,0) -- (2.0,0) node[right] {$x$};
    \draw[->] (0,-1.7) -- (0,1.7) node[above] {$y$};
    \node[below left] at (0,0) {$O$};
    \node[anchor=west] at (0.12,1.0) {I чверть};
    \node[anchor=east] at (-0.12,1.0) {II чверть};
    \node[anchor=east] at (-0.12,-1.0) {III чверть};
    \node[anchor=west] at (0.12,-1.0) {IV чверть};
\end{tikzpicture}\end{nmtfit}
\end{minipage}
\par\penalty-20\vspace{0.25cm}
\end{samepage}

\begin{samepage}
% НМТ 2026, сесія 2.06, завдання №17 --- також у темі: Логарифмічна функція
\zadtask{Установіть відповідність між функцією \mbox{(1--3)} та координатними чвертями \mbox{(А--Д)}, у яких розташований її графік. \nmtyear{2026}

\nmtnobreak
\nopagebreak\nbvspace{0.3cm}
\noindent
\matchingLayout{
\matchHead{Функція}
\matchItem{1}{$y = -x - 1$}
\matchItem{2}{$y = \dfrac{1}{x+2}$}
\matchItem{3}{$y = \log_3 x$}
}{
\matchHead{Чверті}
\matchItem{А}{I і IV}
\matchItem{Б}{I, II і III}
\matchItem{В}{II, III і IV}
\matchItem{Г}{I і III}
\matchItem{Д}{II і IV}
}{
\matchingGrid
}
}
% Відповідь: 1~--~В, 2~--~Б, 3~--~А
\ifshowsolutions\solution{$y=-x-1$ -- спадна пряма з від'ємним зсувом, проходить через II, III і IV чверті -- \textbf{В}; $y=\dfrac{1}{x+2}$ -- гіпербола, зміщена ліворуч, її гілки лежать у I, II і III чвертях -- \textbf{Б}; $y=\log_3 x$ зростає при $x>0$, графік у I і IV чвертях -- \textbf{А}. Відповідь: \textbf{1--В, 2--Б, 3--А}.}\fi
\par\penalty-20
\end{samepage}

\begin{samepage}
\needspace{0.45\textheight}\typeTitle{Найбільше чи найменше значення функції на відрізку}
\noindent\zadnum Установіть відповідність між функцією (1–3) та її найменшим значенням (А–Д) на відрізку $[0; 4]$. \nmtyear{2023}\par
\nopagebreak\nbvspace{0.3cm}

\nmtnobreak
\noindent
\matchingLayout{
\matchHead{Функція}
    \matchItem{1}{$y = x^2 - 2x + 5$}
\matchItem{2}{$y = -0{,}5x + 5$}
\matchItem{3}{$y = 2^x$}
}{
\matchHead{Найменше значення на $[0; 4]$}
    \matchItem{А}{$1$}
\matchItem{Б}{$2$}
\matchItem{В}{$3$}
\matchItem{Г}{$4$}
\matchItem{Д}{$5$}
}{
\matchingGrid
}
\par\penalty-20\vspace{0.25cm}
\end{samepage}

\begin{samepage}
\noindent\zadnum Установіть відповідність між функцією (1–3) та її найбільшим значенням (А–Д) на відрізку $[-1; 3]$. \nmtyear{2023}\par
\nopagebreak\nbvspace{0.3cm}

\nmtnobreak
\noindent
\matchingLayout{
\matchHead{Функція}
    \matchItem{1}{$y = x^2 - 2x$}
\matchItem{2}{$y = -x$}
\matchItem{3}{$y = 5^{-x}$}
}{
\matchHead{Найбільше значення на $[-1; 3]$}
    \matchItem{А}{$1$}
\matchItem{Б}{$2$}
\matchItem{В}{$3$}
\matchItem{Г}{$4$}
\matchItem{Д}{$5$}
}{
\matchingGrid
}
\par\penalty-20\vspace{0.25cm}
\end{samepage}

\begin{samepage}
% НМТ 2026, сесія 4.06, завдання №17
\zadtask{Установіть відповідність між функцією, заданою на проміжку \mbox{(1--3)}, та її найбільшим значенням \mbox{(А--Д)} на цьому проміжку. \nmtyear{2026}

\nmtnobreak
\nopagebreak\nbvspace{0.3cm}
\noindent
\matchingLayout{
\matchHead{Функція}
\matchItem{1}{$y = \sin 4x + 2,\ x \in (-\infty; +\infty)$}
\matchItem{2}{$y = 4 - x^2,\ x \in [0; 2]$}
\matchItem{3}{$y = 4 + \sqrt{x},\ x \in [0; 1]$}
}{
\matchHead{Найбільше значення}
\matchItem{А}{$0$}
\matchItem{Б}{$3$}
\matchItem{В}{$4$}
\matchItem{Г}{$5$}
\matchItem{Д}{$6$}
}{
\matchingGrid
}
}
% Відповідь: 1~--~Б, 2~--~В, 3~--~Г
\ifshowsolutions\solution{1: $\sin4x\leqslant1$, тож найбільше $y=1+2=3$ (Б). 2: $y=4-x^2$ спадає на $[0;2]$, найбільше при $x=0$: $y=4$ (В). 3: $y=4+\sqrt{x}$ зростає на $[0;1]$, найбільше при $x=1$: $y=5$ (Г). Відповідь: \textbf{1~--~Б, 2~--~В, 3~--~Г}.}\fi
\par\penalty-20
\end{samepage}

\begin{samepage}
\needspace{0.45\textheight}\typeTitle{Характеристика функції --- проміжок:\\ $D(y)$, $E(y)$, проміжок монотонності}
% НМТ 2026, сесія 30.05, завдання №17 --- також у темі: Логарифмічна функція
\zadtask{Установіть відповідність між характеристикою функції \mbox{(1--3)} та проміжком \mbox{(А--Д)}, яким є ця характеристика. \nmtyear{2026}

\nmtnobreak
\nopagebreak\nbvspace{0.2cm}
\noindent
\matchingLayout{
\matchHead{Характеристика функції}
\matchItem{1}{область значень функції\\ $y = \sqrt{x+2}$}
\matchItem{2}{область визначення функції \\$y = \lg(x-2)$}
\matchItem{3}{проміжок зростання функції\\ $y = -x^2 + 4x - 7$}
}{
\matchHead{Проміжок}
\matchItem{А}{$(-\infty;\, +\infty)$}
\matchItem{Б}{$(2;\, +\infty)$}
\matchItem{В}{$[0;\, +\infty)$}
\matchItem{Г}{$(-\infty;\, 2]$}
\matchItem{Д}{$(-\infty;\, 2)$}
}{
\matchingGrid
}
}
% Відповідь: 1~--~В, 2~--~Б, 3~--~Г
\ifshowsolutions\solution{\textbf{1}: область значень $y=\sqrt{x+2}$ — це $[0;+\infty)$ (В). \textbf{2}: ОДЗ $\lg(x-2)$: $x-2>0$, тобто $(2;+\infty)$ (Б). \textbf{3}: $y=-x^2+4x-7$ — парабола вітками вниз, вершина $x=2$, зростає на $(-\infty;2]$ (Г). Відповідь: \textbf{1~--~В, 2~--~Б, 3~--~Г}.}\fi
\par\penalty-20
\end{samepage}

\begin{samepage}
\needspace{0.45\textheight}\typeTitle{Кількість спільних точок графіка функції з прямою}
\noindent\zadnum Увідповідніть функцію (1–3) із кількістю (А – Д) спільних точок її графіка з прямою $y=x+3$. \nmtyear{2023}

\nmtnobreak
\nopagebreak\nbvspace{0.3cm}

\nmtnobreak
\noindent
\matchingLayout{
\matchHead{Функція}
    \matchItem{1}{$y=x$}
\matchItem{2}{$y=2^{-x}$}
\matchItem{3}{$y=\dfrac{1}{x}$}
}{
\matchHead{Кількість спільних точок}
    \matchItem{А}{жодної}
\matchItem{Б}{одна}
\matchItem{В}{дві}
\matchItem{Г}{три}
\matchItem{Д}{безліч}
}{
\matchingGrid
}
\par\penalty-20\vspace{0.25cm}
\end{samepage}

\begin{samepage}
\noindent\zadnum Увідповідніть функцію (1–3) із кількістю (А – Д) спільних точок її графіка з прямою $y=1$. \nmtyear{2023}

\nmtnobreak
\nopagebreak\nbvspace{0.3cm}

\nmtnobreak
\noindent
\matchingLayout{
\matchHead{Функція}
    \matchItem{1}{$y=\cos x$}
\matchItem{2}{$y=x^2-2x+2$}
\matchItem{3}{$y=1+\dfrac{2}{x}$}
}{
\matchHead{Кількість спільних точок}
    \matchItem{А}{жодної}
\matchItem{Б}{одна}
\matchItem{В}{дві}
\matchItem{Г}{три}
\matchItem{Д}{безліч}
}{
\matchingGrid
}
\par\penalty-20\vspace{0.25cm}
\end{samepage}

\begin{samepage}
\noindent\zadnum Увідповідніть функцію (1–3) із кількістю (А – Д) спільних точок її графіка з прямою $y=-x$. \nmtyear{2023}

\nmtnobreak
\nopagebreak\nbvspace{0.3cm}

\nmtnobreak
\noindent
\matchingLayout{
\matchHead{Функція}
    \matchItem{1}{$y=-\sqrt{x}$}
\matchItem{2}{$y=\left(\dfrac{1}{2}\right)^x$}
\matchItem{3}{$y=2x+2$}
}{
\matchHead{Кількість спільних точок}
    \matchItem{А}{жодної}
\matchItem{Б}{одна}
\matchItem{В}{дві}
\matchItem{Г}{три}
\matchItem{Д}{безліч}
}{
\matchingGrid
}
\par\penalty-20\vspace{0.25cm}
\end{samepage}

\begin{samepage}
% перенесено з теми 26: завдання про графік або властивості функції
\noindent\zadnum Установіть відповідність між функцією \mbox{(1--3)} та кількістю спільних точок \mbox{(А--Д)} її графіка з прямою $y=x$. \nmtyear{2024}

\nmtnobreak
\nopagebreak\nbvspace{0.3cm}

\nmtnobreak
\matchingLayout{
\matchHead{Функція}
    \matchItem{1}{$y=\dfrac{1}{x}$}
\matchItem{2}{$y=x+3$}
\matchItem{3}{\trigStyle{$y=\operatorname{tg} x$}}

\nmtnobreak
}{
\matchHead{Кількість спільних точок}
    \matchItem{А}{жодної}
\matchItem{Б}{одна}
\matchItem{В}{дві}
\matchItem{Г}{три}
\matchItem{Д}{безліч}
}{
    \matchingGrid
}
\par\penalty-20\vspace{0.25cm}
\end{samepage}

\begin{samepage}
\noindent\zadnum Доберіть до початку речення \mbox{(1--3)} його закінчення \mbox{(А--Д)} так, щоб утворилося правильне твердження. \nmtyear{2024}

\nmtnobreak
\nopagebreak\nbvspace{0.3cm}

\nmtnobreak
\noindent
\matchingLayout{
\matchHead{Початок речення}
    \matchItem{1}{Графік функції $y=x$ не має жодної спільної точки з}
\matchItem{2}{Графік функції $y=3^x$ має лише одну спільну точку з}
\matchItem{3}{Графік рівняння $(x+3)^2+y^2=4$ має дві спільні точки з}
}{
\matchHead{Закінчення речення}
    \matchItem{А}{віссю $x$.}
\matchItem{Б}{віссю $y$.}
\matchItem{В}{прямою $y=x-4$.}
\matchItem{Г}{прямою $y=-4$.}
\matchItem{Д}{прямою $y=-2$.}
}{
\matchingGrid
}
\par\penalty-20\vspace{0.25cm}
\end{samepage}

\begin{samepage}
% === ЗАВДАННЯ 42 (К-сть точок з прямою y=-1) ===
\noindent\zadnum Установіть відповідність між функцією (1–3) та кількістю спільних точок (А – Д) її графіка з прямою $y=-1$. \nmtyear{2024}

\nmtnobreak
\nopagebreak\nbvspace{0.3cm}

\nmtnobreak
\matchingLayout{
\matchHead{Функція}
    \matchItem{1}{$y=\dfrac{2}{x}$}
\matchItem{2}{$y=x^2-4$}
\matchItem{3}{$y=\sin x$}
}{
\matchHead{Кількість спільних точок}
    \matchItem{А}{жодної}
\matchItem{Б}{одна}
\matchItem{В}{дві}
\matchItem{Г}{три}
\matchItem{Д}{більше трьох}
}{
    \matchingGrid
}
\par\penalty-20\vspace{0.25cm}
\end{samepage}

\begin{samepage}
\noindent\zadnum Узгодьте функцію \mbox{(1--3)} із кількістю \mbox{(А--Д)} спільних точок її графіка з прямою $y=-x$. \nmtyear{2025}

\nmtnobreak
\nopagebreak\nbvspace{0.3cm}

\nmtnobreak
\noindent
\matchingLayout{
\matchHead{Функція}
    \matchItem{1}{$y=-x^3$}
\matchItem{2}{$y=2-x$}
\matchItem{3}{$y=2^x$}
}{
\matchHead{Кількість спільних точок}
    \matchItem{А}{жодної}
\matchItem{Б}{одна}
\matchItem{В}{дві}
\matchItem{Г}{три}
\matchItem{Д}{безліч}
}{
\matchingGrid
}
\par\penalty-20\vspace{0.25cm}
\end{samepage}

\begin{samepage}
\needspace{0.45\textheight}\typeTitle{Кількість спільних точок графіка з ламаною, відрізком чи колом на рисунку}
\noindent\zadnum \begin{minipage}[t]{0.55\textwidth}
У прямокутній системі координат на площині зображено замкнену ламану $ABCDA$, де $A(0; -1)$, $B(0; 1)$, $C(2\pi; 1)$, $D(2\pi; -1)$. Увідповідніть функцію (1–3) із кількістю (А – Д) спільних точок її графіка з ламаною $ABCDA$. \nmtyear{2023}
\end{minipage}
\hfill
\begin{minipage}[t]{0.4\textwidth}
    \nopagebreak\nbvspace{-0.3cm}
    \begin{flushright}
    \begin{nmtfit}\begin{tikzpicture}[scale=1.31, x=0.5cm, y=1cm]
        % Сітка не потрібна, лише осі та прямокутник
        \draw[->, thick] (-1,0) -- (7.5,0) node[below] {$x$};
        \draw[->, thick] (0,-1.5) -- (0,1.8) node[left] {$y$};

        % Ламана
        \coordinate (A) at (0, -1);
        \coordinate (B) at (0, 1);
        \coordinate (C) at (6.28, 1); % 2pi approx 6.28
        \coordinate (D) at (6.28, -1);

        \draw[thick] (A) -- (B) -- (C) -- (D) -- cycle;

        % Підписи точок
        \node[below right] at (A) {$A$};
        \node[above right] at (B) {$B$};
        \node[above right] at (C) {$C$};
        \node[below right] at (D) {$D$};

        % Підписи осей
        \node[below left] at (0,0) {$0$};
        \node[left] at (0,1) {$1$};
        \node[left] at (0,-1) {$-1$};
        \node[below] at (6.28, 0) {$2\pi$};
        \draw (1.57, 0.1) -- (1.57, -0.1) node[below] {$\frac{\pi}{2}$};
        \draw (3.14, 0.1) -- (3.14, -0.1) node[below] {$\pi$};
        \draw (4.71, 0.1) -- (4.71, -0.1) node[below] {$\frac{3\pi}{2}$};

    \end{tikzpicture}\end{nmtfit}
    \end{flushright}
\end{minipage}

\nmtnobreak
\nopagebreak\nbvspace{0.3cm}

\nmtnobreak
\matchingLayout{
\matchHead{Функція}
    \matchItem{1}{$y=x^3$}
\matchItem{2}{$y=\cos x$}
\matchItem{3}{$y=x+1$}
}{
\matchHead{К-сть спільних точок}
    \matchItem{А}{жодної}
\matchItem{Б}{одна}
\matchItem{В}{дві}
\matchItem{Г}{три}
\matchItem{Д}{чотири}
}{
    \matchingGrid
}
\par\penalty-20\vspace{0.25cm}
\end{samepage}

\begin{samepage}
\noindent\zadnum \begin{minipage}[t]{0.55\textwidth}
У прямокутній системі координат на площині зображено коло $x^2+y^2=4$. Установіть відповідність між функцією (1–3) та кількістю (А – Д) спільних точок, які має графік цієї функції із заданим колом. \nmtyear{2023}
\end{minipage}
\hfill
\begin{minipage}[t]{0.4\textwidth}
    \nopagebreak\nbvspace{-0.3cm}
    \begin{flushright}
    \begin{nmtfit}\begin{tikzpicture}[scale=0.74]
        % Сітка та осі

        \draw[->, thick] (-3,0) -- (3,0) node[below] {$x$};
        \draw[->, thick] (0,-3) -- (0,3) node[left] {$y$};

        % Коло
        \draw[thick] (0,0) circle (2cm);

        % Підписи
        \node[below left] at (0,0) {$0$};
        \node[above right] at (1.4, 1.4) {$x^2+y^2=4$};
    \end{tikzpicture}\end{nmtfit}
    \end{flushright}
\end{minipage}

\nmtnobreak
\nopagebreak\nbvspace{0.3cm}

\nmtnobreak
\noindent
\matchingLayout{
\matchHead{Функція}
    \matchItem{1}{$y=x+4$}
\matchItem{2}{$y=x^2-2$}
\matchItem{3}{$y=4^x$}
}{
\matchHead{Кількість спільних точок}
    \matchItem{А}{жодної}
\matchItem{Б}{лише одна}
\matchItem{В}{лише дві}
\matchItem{Г}{лише три}
\matchItem{Д}{більше за три}
}{
\matchingGrid
}
\par\penalty-20\vspace{0.25cm}
\end{samepage}

\begin{samepage}
\noindent\zadnum \begin{minipage}[t]{0.55\textwidth}
У прямокутній системі координат на площині зображено ламану $ABCDA$ (див. рисунок). Увідповідніть функцію (1–3) із кількістю (А – Д) спільних точок її графіка з ламаною $ABCDA$. \nmtyear{2023}
\end{minipage}
\hfill
\begin{minipage}[t]{0.4\textwidth}
    \nopagebreak\nbvspace{-0.3cm}
    \begin{flushright}
    \begin{nmtfit}\begin{tikzpicture}[scale=0.5]
    % одиниця осі x -- 2 клітинки (за оригіналом): π/2 ≈ 3,14 клітинки; одиниця осі y -- 2 клітинки
    \draw[nmt grid, step=1] (-5,-4) grid (6,5);
    \draw[->] (-5,0) -- (6,0) node[below left] {$x$};
    \draw[->] (0,-4) -- (0,5) node[below left] {$y$};
    \node[below left] at (0,0) {$0$};
    \node[above left] at (0,2) {$1$};
    \node[below left] at (0,-2) {$-1$};

    \coordinate (A) at (-3.14,-2);
    \coordinate (B) at (-3.14,2);
    \coordinate (C) at (3.14,2);
    \coordinate (D) at (3.14,-2);
    \draw[nmt curve] (A) -- (B) -- (C) -- (D) -- cycle;
    \node[nmt dot] at (A) {};
    \node[nmt dot] at (B) {};
    \node[nmt dot] at (C) {};
    \node[nmt dot] at (D) {};
    \node[below left] at (A) {$A$};
    \node[above left] at (B) {$B$};
    \node[above right] at (C) {$C$};
    \node[below right] at (D) {$D$};
    \node[below left, xshift=-2pt] at (-3.14,0) {$-\dfrac{\pi}{2}$};
    \node[below right, xshift=2pt] at (3.14,0) {$\dfrac{\pi}{2}$};
\end{tikzpicture}\end{nmtfit}
    \end{flushright}
\end{minipage}

\nmtnobreak
\nopagebreak\nbvspace{0.3cm}

\nmtnobreak
\noindent
\matchingLayout{
\matchHead{Функція}
    \matchItem{1}{$y=\dfrac{\pi}{2}$}
\matchItem{2}{$y=\cos x$}
\matchItem{3}{$y=x^2+1$}
}{
\matchHead{Кількість спільних точок}
    \matchItem{А}{жодної}
\matchItem{Б}{одна}
\matchItem{В}{дві}
\matchItem{Г}{три}
\matchItem{Д}{безліч}
}{
\matchingGrid
}
\par\penalty-20\vspace{0.25cm}
\end{samepage}

\begin{samepage}
% перенесено з теми 26: завдання про графік або властивості функції
\noindent\zadnum \begin{minipage}[t]{0.55\textwidth}
У прямокутній системі координат на площині зображено ламану $ABC$, $A(-2; 0)$, $B(0; 1)$, $C(2; 1)$ (див. рисунок). Установіть відповідність між функцією \mbox{(1--3)} та кількістю спільних точок \mbox{(А--Д)} її графіка з ламаною $ABC$.
 \nmtyear{2024}
\end{minipage}
\hfill
\begin{minipage}[t]{0.4\textwidth}
    \nopagebreak\nbvspace{-0.3cm}
    \begin{flushright}
    \begin{nmtfit}\begin{tikzpicture}[scale=0.5]
        \draw[nmt grid, step=1cm] (-2.5,-0.5) grid (2.5,2.5);
        \draw[->, thick] (-2.5,0) -- (2.5,0) node[below] {$x$};
        \draw[->, thick] (0,-0.5) -- (0,2.5) node[left] {$y$};
        \node[below left] at (0,0) {$0$};

        \draw[thick] (-2,0) -- (0,1) -- (2,1);

        \path (-2,0) node[nmt dot] {} node[above left] {$A$};
        \path (0,1) node[nmt dot] {} node[above right] {$B$};
        \path (2,1) node[nmt dot] {} node[above right] {$C$};

        \node[below] at (-2,0) {$-2$};
        \node[below] at (1,0) {$1$};
        \node[below] at (2,0) {$2$};
        \node[left] at (0,1) {$1$};
    \end{tikzpicture}\end{nmtfit}
    \end{flushright}
\end{minipage}

\nmtnobreak
\nopagebreak\nbvspace{0.3cm}

\nmtnobreak
\matchingLayout{
\matchHead{Функція}
    \matchItem{1}{$y=2-x^2$}
\matchItem{2}{\trigStyle{$y=\sin x$}}
\matchItem{3}{$y=\log_5 x$}
}{
\matchHead{Кількість спільних точок}
    \matchItem{А}{жодної}
\matchItem{Б}{одна}
\matchItem{В}{дві}
\matchItem{Г}{три}
\matchItem{Д}{більше трьох}
}{
    \matchingGrid
}
\par\penalty-20\vspace{0.25cm}
\end{samepage}

\begin{samepage}
% === ЗАВДАННЯ 49 (Ламана ABCA і функції) ===
\noindent\zadnum \begin{minipage}[t]{0.55\textwidth}
У прямокутній декартовій системі координат на площині зображено замкнену ламану $ABCA$, де $A(-1; 0)$, $B(0; 1)$, $C(1; 0)$. Узгодьте функцію (1–3) з кількістю (А – Д) спільних точок її графіка та ламаної $ABCA$. \nmtyear{2024}
\end{minipage}
\hfill
\begin{minipage}[t]{0.4\textwidth}
    \nopagebreak\nbvspace{-0.3cm}
    \begin{flushright}
    \begin{nmtfit}\begin{tikzpicture}[scale=2.04]
        % Осі
        \draw[->, thick] (-1.5,0) -- (1.5,0) node[below] {$x$};
        \draw[->, thick] (0,-0.5) -- (0,1.5) node[left] {$y$};

        % Ламана (зеленим кольором, як на рисунку)
        \draw[thick, nmtGreen] (-1, 0) -- (0, 1) -- (1, 0) -- cycle;

        % Точки
        \path (-1,0) node[nmt dot] {} node[above left] {$A$};
        \path (0,1) node[nmt dot] {} node[above right] {$B$};
        \path (1,0) node[nmt dot] {} node[above right] {$C$};

        % Підписи координат
        \node[below] at (-1,0) {$-1$};
        \node[below left] at (0,0) {$0$};
        \node[below] at (1,0) {$1$};
        \node[left] at (0,1) {$1$};
    \end{tikzpicture}\end{nmtfit}
    \end{flushright}
\end{minipage}

\nmtnobreak
\nopagebreak\nbvspace{0.3cm}

\nmtnobreak
\matchingLayout{
\matchHead{Функція}
    \matchItem{1}{$y=0$}
\matchItem{2}{$y=1-x^2$}
\matchItem{3}{$y=\cos x$}
}{
\matchHead{Кількість спільних точок}
    \matchItem{А}{жодної}
\matchItem{Б}{лише одна}
\matchItem{В}{лише дві}
\matchItem{Г}{лише три}
\matchItem{Д}{безліч}
}{
    \matchingGrid
}
\par\penalty-20\vspace{0.25cm}
\end{samepage}

\begin{samepage}
% НМТ 2026, сесія 11.06, завдання №17
\zadtask{У прямокутній системі координат на площині $xy$ зображено відрізок $AB$ (див. рисунок). Увідповідніть функцію \mbox{(1--3)} із кількістю \mbox{(А--Д)} спільних точок її графіка з відрізком $AB$. \nmtyear{2026}\par
\nopagebreak\nbvspace{0.3cm}
\begin{center}
\begin{nmtfit}\begin{tikzpicture}[scale=1.1]
    \draw[->] (-2,0) -- (2,0) node[right] {$x$};
    \draw[->] (0,-1) -- (0,3) node[above] {$y$};
    \draw[thin] (-1,0.08) -- (-1,-0.08); \draw[thin] (1,0.08) -- (1,-0.08);
    \draw[thin] (-0.08,1) -- (0.08,1); \draw[thin] (-0.08,2) -- (0.08,2);
    \node[below left] at (0,0) {$0$};
    \node[below] at (-1,0) {$-1$};
    \node[below] at (1,0) {$1$};
    \node[right, xshift=2pt] at (0,1) {$1$};
    \node[right, xshift=2pt] at (0,2) {$2$};
    \draw[dashed, thin, nmtGreen] (-1,2) -- (-1,0);
    \draw[dashed, thin, nmtGreen] (-1,2) -- (0,2);
    \draw[thick] (-1,2) -- (0,1);
    \node[nmt dot] at (-1,2) {};
    \node[nmt dot] at (0,1) {};
    \node[left] at (-1,2) {$B$};
    \node[below left] at (0,1) {$A$};
\end{tikzpicture}\end{nmtfit}
\end{center}
\nopagebreak\nbvspace{0.2cm}
\noindent
\matchingLayout{
\matchHead{Функція}
\matchItem{1}{$y = 0{,}5^{\,x}$}
\matchItem{2}{$y = \cos x$}
\matchItem{3}{$y = 1 - x$}
}{
\matchHead{Кількість спільних точок}
\matchItem{А}{жодної}
\matchItem{Б}{лише одна}
\matchItem{В}{лише дві}
\matchItem{Г}{лише три}
\matchItem{Д}{безліч}
}{
\matchingGrid
}
}
\par\penalty-20
\end{samepage}

\begin{samepage}
\needspace{0.45\textheight}\typeTitle{Графік функції на рисунку: значення, нулі, точки екстремуму, зсуви}
\noindent\zadnum \begin{minipage}[t]{0.55\textwidth}
На рисунку зображено графік функції $y=f(x)$, визначеної на проміжку $[-4; 5]$. Установіть відповідність між початком речення \mbox{(1--3)} та його закінченням \mbox{(А--Д)} так, щоб утворилося правильне твердження. \nmtyear{2024}
\end{minipage}
\hfill
\begin{minipage}[t]{0.4\textwidth}
\nbvspace{0pt}
\begin{flushright}
\begin{nmtfit}\begin{tikzpicture}[scale=0.5]
        % Сітка розширена вниз до -3.5
        \draw[nmt grid, step=1cm] (-4.5,-3.5) grid (5.5,4.5);
        \draw[->, thick] (-4.5,0) -- (5.5,0) node[below] {$x$};
        \draw[->, thick] (0,-3.5) -- (0,4.5) node[left] {$y$};

        \node[below left] at (0,0) {$0$};
        \node[below] at (1,0) {$1$};
        \node[left] at (0,1) {$1$};
        \node[below] at (5,0) {$5$};
        \node[below] at (-4,0) {$-4$};

        % Графік строго через точки: (-4, -3), (-1.2, 0), (0, 2.4), (2.5, 4), (5, 2)
        \draw[nmt curve] plot [smooth, tension=0.6] coordinates {(-4, -3) (-1.2, 0) (0, 2.4) (2.5, 4) (5, 2)};

        \path (-4,-3) node[nmt dot] {};
        \path (5,2) node[nmt dot] {};
        \node[above, fill=white, inner sep=1.5pt] at (4.2, 4.5) {$y=f(x)$};
    \end{tikzpicture}\end{nmtfit}
\end{flushright}
\end{minipage}

\nmtnobreak
\nopagebreak\nbvspace{0.3cm}

\nmtnobreak
\noindent
\matchingLayout{
\matchHead{Початок речення}
\matchItem{1}{Нуль функції належить проміжку}
\matchItem{2}{Точка максимуму функції належить проміжку}
\matchItem{3}{Абсциса точки перетину графіка функції з графіком функції $y = \log_{\frac{1}{3}} x$ належить проміжку}
}{
\matchHead{Закінчення речення}
    \matchItem{А}{$(-4; -2]$.}
\matchItem{Б}{$(-2; 0]$.}
\matchItem{В}{$(0; 1]$.}
\matchItem{Г}{$(1; 3]$.}
\matchItem{Д}{$(3; 5]$.}
}{
\matchingGrid
}
\par\penalty-20\vspace{0.25cm}
\end{samepage}

\begin{samepage}
% НМТ 2026, сесія 10.06, завдання №17 --- основна тема: Похідна
\zadtask{На рисунку зображено графік функції $y = f(x)$, визначеної на проміжку $[-2;\, 4]$. Для будь-якого значення $x$ із проміжку $(-2;\, 4)$ існує похідна $f'(x)$. До кожного початку \mbox{(1--3)} речення доберіть його закінчення \mbox{(А--Д)} так, щоб утворилося правильне твердження. \nmtyear{2026}\par
\nopagebreak\nbvspace{0.3cm}
\begin{center}
\begin{nmtfit}\begin{tikzpicture}[scale=0.5]
    \draw[nmt grid, step=1] (-3,-4) grid (6,6);
    \draw[->, thick] (-3,0) -- (6,0) node[below left] {$x$};
    \draw[->, thick] (0,-4) -- (0,6) node[left] {$y$};
    \draw[thin] (-2,0.12) -- (-2,-0.12); \draw[thin] (1,0.12) -- (1,-0.12); \draw[thin] (4,0.12) -- (4,-0.12); \draw[thin] (0.12,1) -- (-0.12,1); \draw[thin] (0.12,5) -- (-0.12,5); \draw[thin] (0.12,-3) -- (-0.12,-3);
    \node[below left] at (0,0) {$0$};
    \node[below left] at (-2,0) {$-2$};
    \node[below] at (1,0) {$1$};
    \node[below] at (4,0) {$4$};
    \node[right] at (0,1) {$1$};
    \node[right] at (0,5) {$5$};
    \node[left] at (0,-3) {$-3$};
    \draw[nmt curve] (-2,-3) .. controls (-1.67,-1.6) and (-1.33,-1) .. (-1,0) .. controls (-0.67,1) and (-0.33,1.4) .. (0,2) .. controls (0.33,2.6) and (0.65,3) .. (1,3) .. controls (1.4,3) and (1.7,1) .. (2,1) .. controls (2.8,1) and (3.2,2.44) .. (4,5);
    \node[nmt dot] at (-2,-3) {};
    \node[nmt dot] at (4,5) {};
    \node[fill=white, inner sep=1.5pt] at (4.7,1.6) {$y=f(x)$};
\end{tikzpicture}\end{nmtfit}
\end{center}
\nopagebreak\nbvspace{0.2cm}
\noindent
\matchingLayout{
\matchHead{Початок речення}
\matchItem{1}{Найбільше значення функції $y=f(x)$ на проміжку $[-2;\, 4]$ дорівнює}
\matchItem{2}{Функція набуває від'ємного значення, якщо $x =$}
\matchItem{3}{$f'(x) = 0$, якщо $x =$}
}{
\matchHead{Закінчення речення}
\matchItem{А}{$-2$.}
\matchItem{Б}{$-1$.}
\matchItem{В}{$2$.}
\matchItem{Г}{$4$.}
\matchItem{Д}{$5$.}
}{
\matchingGrid
}
}
\par\penalty-20
\end{samepage}

\begin{samepage}
\needspace{0.45\textheight}\typeTitle{Графік із кількох частин: проміжок --- формула}
\noindent\zadnum \begin{minipage}[t]{0.55\textwidth}
На рисунку зображено графік функції $y=f(x)$, визначеної на проміжку $[-4; 6]$. Узгодьте проміжок \mbox{(1--3)} із функцією \mbox{(А--Д)}, що описує графік функції $y=f(x)$ на цьому проміжку. \nmtyear{2025}
\end{minipage}
\hfill
\begin{minipage}[t]{0.4\textwidth}
    \nopagebreak\nbvspace{-0.3cm}
    \begin{flushright}
    \begin{nmtfit}\begin{tikzpicture}[scale=0.5]
        \draw[nmt grid, step=1] (-4.5,-0.5) grid (6.5,8.5);
        \draw[->, thick] (-4.5,0) -- (6.5,0) node[below] {$x$};
        \draw[->, thick] (0,-0.5) -- (0,8.5) node[left] {$y$};
        \node[below left] at (0,0) {$0$};
        \node[below] at (1,0) {$1$};
        \node[left] at (0,1) {$1$};
        \node[below] at (6,0) {$6$};
        \node[below] at (-4,0) {$-4$};
        % [-4; 0]: y = sqrt(-x); (0; 2]: y = x^3; (2; 6]: y = 12 - 2x
        \draw[nmt curve] plot [smooth, domain=-4:0] (\x, {sqrt(-\x)});
        \draw[nmt curve] plot [smooth, domain=0:2] (\x, {\x*\x*\x});
        \draw[nmt curve] (2, 8) -- (6, 0);
        \node[nmt dot] at (-4, 2) {};
        \node[nmt dot] at (6, 0) {};
        \node[anchor=west, fill=white, inner sep=1.5pt] at (3.7, 6.7) {$y=f(x)$};
\end{tikzpicture}\end{nmtfit}
    \end{flushright}
\end{minipage}

\nmtnobreak
\nopagebreak\nbvspace{0.3cm}

\nmtnobreak
\noindent
\matchingLayout{
\matchHead{Проміжок}
    \matchItem{1}{$[-4; 0]$}
\matchItem{2}{$(0; 2]$}
\matchItem{3}{$(2; 6]$}
}{
\matchHead{Функція}
    \matchItem{А}{$y=x^3$}
\matchItem{Б}{$y=12-2x$}
\matchItem{В}{$y=-\sqrt{x}$}
\matchItem{Г}{$y=2^x$}
\matchItem{Д}{$y=\sqrt{-x}$}
}{
\matchingGrid
}
\par\penalty-20\vspace{0.25cm}
\end{samepage}

\begin{samepage}
% НМТ 2026, сесія 18.06, завдання №17 --- також у темі: Показникова функція і показникові вирази
\zadtask{На рисунку зображено графік функції $y = f(x)$, визначеної на проміжку $(-\infty;\, +\infty)$. Узгодьте проміжок \mbox{(1--3)} із функцією \mbox{(А--Д)}, що описує графік функції $y = f(x)$ на цьому проміжку. \nmtyear{2026}\par
\nopagebreak\nbvspace{0.3cm}
\begin{center}
\begin{nmtfit}\begin{tikzpicture}[scale=0.5]
    \draw[nmt grid, step=1] (-6,-2) grid (8,4);
    \draw[->] (-6.3,0) -- (8.4,0) node[right] {$x$};
    \draw[->] (0,-2.3) -- (0,4.4) node[above] {$y$};
    \node[below left] at (0,0) {$0$};
    \draw[thin] (1,0.12) -- (1,-0.12);
    \node[below] at (1,0) {$1$};
    \draw[thin] (0.12,1) -- (-0.12,1);
    \node[right] at (0,1) {$1$};
    \draw[nmt curve] plot[domain=-6:-1, samples=50] (\x, {-2/\x}) -- (0,0);
    \draw[nmt curve] plot[domain=0:7.8540, samples=100] (\x, {sin(deg(\x))});
    \node[right, fill=white, inner sep=1.5pt] at (1.5,2.4) {$y=f(x)$};
\end{tikzpicture}\end{nmtfit}
\end{center}
\nopagebreak\nbvspace{0.2cm}
\noindent
\matchingLayout{
\matchHead{Проміжок}
\matchItem{1}{$(-\infty;\, -1)$}
\matchItem{2}{$[-1;\, 0]$}
\matchItem{3}{$(0;\, +\infty)$}
}{
\matchHead{Функція}
\matchItem{А}{$y = \cos x$}
\matchItem{Б}{$y = 2^x$}
\matchItem{В}{$y = -\dfrac{2}{x}$}
\matchItem{Г}{$y = \sin x$}
\matchItem{Д}{$y = -2x$}
}{
\matchingGrid
}
}
% Відповідь: 1--В,\ \ 2--Д,\ \ 3--Г
\par\penalty-20
\end{samepage}

\begin{samepage}
\needspace{0.45\textheight}\typeTitle{Три рисунки графіків --- твердження про функцію}
\noindent\zadnum На кожному з рисунків \mbox{(1--3)} зображено графік функції $y=f(x)$, визначеної на проміжку $[-3; 3]$. Узгодьте рисунок \mbox{(1--3)} з твердженням \mbox{(А--Д)} щодо функції $y=f(x)$, графік якої зображено на цьому рисунку. \nmtyear{2024}

\nmtnobreak
\nopagebreak\nbvspace{0.3cm}

\nmtnobreak
% Блок рисунків
\noindent
\begin{minipage}[t]{0.32\textwidth}\nbvspace{0pt}
    \centering
    Рис. 1 \\
    \begin{nmtfit}\begin{tikzpicture}[scale=0.3]
    \draw[nmt grid, step=1] (-5,-5) grid (5,3);
    \draw[->] (-5,0) -- (5,0) node[above left] {$x$};
    \draw[->] (0,-5) -- (0,3) node[left] {$y$};
    \node[below left] at (0,0) {$0$};
    \node[below] at (1,0) {$1$};
    \node[left] at (0,1) {$1$};
    \node[left] at (-3,-1) {$-3$};
    \node[right] at (3,-1) {$3$};

    \draw[nmt curve] (-3,-1) arc[start angle=180, end angle=360, radius=3];
    \node[nmt dot] at (-3,-1) {};
    \node[nmt dot] at (3,-1) {};
    \node[anchor=north east, inner sep=1.5pt] at (5,-5) {$y=f(x)$};
\end{tikzpicture}\end{nmtfit}
\end{minipage}
\hfill
\begin{minipage}[t]{0.32\textwidth}\nbvspace{0pt}
    \centering
    Рис. 2 \\
    \begin{nmtfit}\begin{tikzpicture}[scale=0.3]
    \draw[nmt grid, step=1] (-5,-3) grid (5,5);
    \draw[->] (-5,0) -- (5,0) node[above left] {$x$};
    \draw[->] (0,-3) -- (0,5) node[left] {$y$};
    \node[below left] at (0,0) {$0$};
    \node[above right] at (1,0) {$1$};
    \node[right] at (0,1) {$1$};
    \node[below] at (-3,0) {$-3$};
    \node[below] at (3,0) {$3$};

    \draw[nmt curve] (-3,0) .. controls (-2.8,2.6) and (-2.1,4) .. (-1.5,4)
        .. controls (-1.1,4) and (-0.5,1) .. (0,0)
        .. controls (0.4,-0.9) and (1.4,-1.5) .. (3,-1.75);
    \node[nmt dot] at (-3,0) {};
    \node[nmt dot] at (3,-1.75) {};
    \node[anchor=north east, inner sep=1.5pt, yshift=-3pt] at (5,-3) {$y=f(x)$};
\end{tikzpicture}\end{nmtfit}
\end{minipage}
\hfill
\begin{minipage}[t]{0.32\textwidth}\nbvspace{0pt}
    \centering
    Рис. 3 \\
    \begin{nmtfit}\begin{tikzpicture}[scale=0.3]
    \draw[nmt grid, step=1] (-5,-3) grid (5,5);
    \draw[->] (-5,0) -- (5,0) node[above left] {$x$};
    \draw[->] (0,-3) -- (0,5) node[left] {$y$};
    \node[below left] at (0,0) {$0$};
    \node[above right] at (1,0) {$1$};
    \node[right] at (0,1) {$1$};
    \node[below] at (-3,0) {$-3$};
    \node[below] at (3,0) {$3$};

    \draw[nmt curve] (-3,5) .. controls (-2.9,3.6) and (-2.4,2.4) .. (-2,2) -- (3,-3);
    \node[nmt dot] at (-3,5) {};
    \node[nmt dot] at (3,-3) {};
    \node[anchor=north east, inner sep=1.5pt, yshift=-3pt] at (5,-3) {$y=f(x)$};
\end{tikzpicture}\end{nmtfit}
\end{minipage}

\nmtnobreak
\nopagebreak\nbvspace{0.3cm}

\nmtnobreak
\noindent
\matchingLayout{
\matchHead{Рисунок}
\matchItem{1}{Рис. 1}
\matchItem{2}{Рис. 2}
\matchItem{3}{Рис. 3}
}{
\matchHead{Твердження}
    \matchItem{А}{графік функції $f$ є фрагментом кола $x^2 + (y-1)^2 = 9$}
\matchItem{Б}{графік функції $f$ є фрагментом кола $x^2 + (y+1)^2 = 9$}
\matchItem{В}{функція $f$ зростає на області визначення}
\matchItem{Г}{графік функції $y=f(x)+2$ розташований лише в I і II чвертях}
\matchItem{Д}{графік функції $f$ має лише одну спільну точку з графіком функції $y=2^x$}
}{
\matchingGrid
}
\par\penalty-20\vspace{0.25cm}
\end{samepage}

\begin{samepage}
\noindent\zadnum На кожному з рисунків \mbox{(1--3)} зображено графік функції $y=f(x)$, визначеної на проміжку $[-3; 3]$. Узгодьте рисунок \mbox{(1--3)} з твердженням \mbox{(А--Д)} щодо функції $y=f(x)$, графік якої зображено на цьому рисунку. \nmtyear{2025}

\nmtnobreak
\nopagebreak\nbvspace{0.3cm}

\nmtnobreak
% Блок рисунків
\noindent
\begin{minipage}[t]{0.32\textwidth}\nbvspace{0pt}
    \centering
    Рис. 1 \\
    \begin{nmtfit}\begin{tikzpicture}[scale=0.3]
        \draw[nmt grid, step=1] (-5,-3) grid (5,5);
        \draw[->] (-5,0) -- (5.3,0) node[below left] {$x$};
        \draw[->] (0,-3) -- (0,5.3) node[left] {$y$};
        \node[below left] at (0,0) {\small $0$};
        \node[above left] at (0,1) {\small $1$};
        \node[below] at (1,0) {\small $1$};
        \node[below] at (3,0) {\small $3$};
        \node[below] at (-3,0) {\small $-3$};
        \draw[nmt curve] plot[domain=0:3, samples=40] (\x, {1+1.1547*sqrt(\x)});
        \draw[nmt curve] plot[domain=-3:0, samples=40] (\x, {1+1.1547*sqrt(-\x)});
        \path (-3,3) node[nmt dot] {};
        \path (3,3) node[nmt dot] {};
        \node[fill=white, inner sep=1pt] at (3.0,3.95) {\small $y=f(x)$};
\end{tikzpicture}\end{nmtfit}
\end{minipage}
\hfill
\begin{minipage}[t]{0.32\textwidth}\nbvspace{0pt}
    \centering
    Рис. 2 \\
    \begin{nmtfit}\begin{tikzpicture}[scale=0.3]
        \draw[nmt grid, step=1] (-5,-3) grid (5,5);
        \draw[->] (-5,0) -- (5.3,0) node[below left] {$x$};
        \draw[->] (0,-3) -- (0,5.3) node[left] {$y$};
        \node[below left] at (0,0) {\small $0$};
        \node[above left] at (0,1) {\small $1$};
        \node[below] at (1,0) {\small $1$};
        \node[below] at (3,0) {\small $3$};
        \node[below] at (-3,0) {\small $-3$};
        \draw[nmt curve] (-3,4) -- (2,0) -- (3,-2);
        \path (-3,4) node[nmt dot] {};
        \path (3,-2) node[nmt dot] {};
        \node[fill=white, inner sep=1pt] at (3.0,3.95) {\small $y=f(x)$};
\end{tikzpicture}\end{nmtfit}
\end{minipage}
\hfill
\begin{minipage}[t]{0.32\textwidth}\nbvspace{0pt}
    \centering
    Рис. 3 \\
    \begin{nmtfit}\begin{tikzpicture}[scale=0.3]
        \draw[nmt grid, step=1] (-5,-3) grid (5,5);
        \draw[->] (-5,0) -- (5.3,0) node[below left] {$x$};
        \draw[->] (0,-3) -- (0,5.3) node[left] {$y$};
        \node[below left] at (0,0) {\small $0$};
        \node[above left] at (0,1) {\small $1$};
        \node[below] at (1,0) {\small $1$};
        \node[below] at (3,0) {\small $3$};
        \node[below] at (-3,0) {\small $-3$};
        \draw[nmt curve] (-3,2) .. controls (-2.667,0.333) and (-2.333,-1) .. (-2,-1) .. controls (-1.667,-1) and (-1.333,-0.333) .. (-1,0) .. controls (-0.667,0.333) and (-0.333,0.667) .. (0,1) .. controls (0.333,1.333) and (0.667,1.769) .. (1,2) .. controls (1.667,2.462) and (2.333,2.889) .. (3,3);
        \path (-3,2) node[nmt dot] {};
        \path (3,3) node[nmt dot] {};
        \node[fill=white, inner sep=1pt] at (3.0,3.95) {\small $y=f(x)$};
\end{tikzpicture}\end{nmtfit}
\end{minipage}

\nmtnobreak
\nopagebreak\nbvspace{0.3cm}

\nmtnobreak
\noindent
\matchingLayout{
\matchHead{Рисунок}
\matchItem{1}{Рис. 1}
\matchItem{2}{Рис. 2}
\matchItem{3}{Рис. 3}
}{
\matchHead{Твердження}
    \matchItem{А}{графік функції $f$ має лише одну спільну точку з графіком кола $(x-3)^2 + (y+2)^2 = 1$}
\matchItem{Б}{похідна функції $f$ є додатною на проміжку $[-3; 3]$}
\matchItem{В}{функція $f$ є непарною}
\matchItem{Г}{значення функції $f$ на кінцях проміжку $[-3; 3]$ є рівними}
\matchItem{Д}{функція $f$ має два нулі}
}{
\matchingGrid
}
\par\penalty-20\vspace{0.25cm}
\end{samepage}

\begin{samepage}
\needspace{0.45\textheight}\typeTitle{Прямі на рисунках --- твердження:\\ рівняння, кутовий коефіцієнт, спільні точки}
\noindent\zadnum Установіть відповідність між твердженням \mbox{(1--3)} та прямою, зображеною на рисунку \mbox{(А--Д)}, для якої це твердження є правильним. \nmtyear{2024}

\nmtnobreak
\nopagebreak\nbvspace{0.3cm}

\nmtnobreak
\noindent
\begin{minipage}[t]{0.60\textwidth}
\matchHead{Твердження}
    \matchItem{1}{не має спільних точок з функцією $y=\left(\dfrac{2}{3}\right)^x$}
\matchItem{2}{є графіком функції $y=x-2$}
\matchItem{3}{кутовий коефіцієнт прямої дорівнює $0$}
\end{minipage}%
\hfill
\begin{minipage}[t]{0.30\textwidth}
    \nopagebreak\nbvspace{0.5cm}
    \matchingGrid
\end{minipage}

\nmtnobreak
\nopagebreak\nbvspace{0.5cm}
\matchHead{Пряма}

\nmtnobreak
% Таблиця з графіками
\begin{center}
\begingroup
\setlength{\tabcolsep}{3pt}
\begin{tabular}{|*{5}{c|}}
\hline
\textbf{А} & \textbf{Б} & \textbf{В} & \textbf{Г} & \textbf{Д} \\
\hline
\begin{nmtfit}\begin{tikzpicture}[scale=0.33, thick]
    \useasboundingbox (-3.8,-3.4) rectangle (4.6,3.5);
    \draw[->] (-3.4,0) -- (4.2,0) node[below] {$x$};
    \draw[->] (0,-3.2) -- (0,3.0) node[left] {$y$};
    \node[below left] at (0,0) {\small $0$};
    \node[below right] at (2,0) {\small $2$};
    \node[below right] at (0,-2) {\small $-2$};
    \draw (-1.2,-3.2) -- (3.3,1.3);
\end{tikzpicture}\end{nmtfit} &
\begin{nmtfit}\begin{tikzpicture}[scale=0.33, thick]
    \useasboundingbox (-3.8,-3.4) rectangle (4.6,3.5);
    \draw[->] (-3.4,0) -- (4.2,0) node[below] {$x$};
    \draw[->] (0,-3.2) -- (0,3.0) node[left] {$y$};
    \coordinate (O) at (0,0); \coordinate (P) at (0,1); \coordinate (Q) at (1,1);
    \node[below left] at (0,0) {\small $0$};
    \node[above left] at (0,1) {\small $1$};
    \pic[draw, angle radius=0.17cm] {right angle = Q--P--O};
    \draw (-3.4,1) -- (3.4,1);
\end{tikzpicture}\end{nmtfit} &
\begin{nmtfit}\begin{tikzpicture}[scale=0.33, thick]
    \useasboundingbox (-3.8,-3.4) rectangle (4.6,3.5);
    \draw[->] (-3.4,0) -- (4.2,0) node[below] {$x$};
    \draw[->] (0,-3.2) -- (0,3.0) node[right] {$y$};
    \coordinate (O) at (0,0); \coordinate (P) at (-1,0); \coordinate (Q) at (-1,1);
    \node[below right] at (0,0) {\small $0$};
    \node[below left, xshift=-2pt] at (-1,0) {\small $-1$};
    \pic[draw, angle radius=0.17cm] {right angle = O--P--Q};
    \draw (-1,-3.2) -- (-1,3.0);
\end{tikzpicture}\end{nmtfit} &
\begin{nmtfit}\begin{tikzpicture}[scale=0.33, thick]
    \useasboundingbox (-3.8,-3.4) rectangle (4.6,3.5);
    \draw[->] (-3.4,0) -- (4.2,0) node[below] {$x$};
    \draw[->] (0,-3.2) -- (0,3.0) node[left] {$y$};
    \node[below right] at (0,0) {\small $0$};
    \node[below left] at (-2,0) {\small $-2$};
    \node[below left] at (0,-2) {\small $-2$};
    \draw (-3.2,1.2) -- (1.2,-3.2);
\end{tikzpicture}\end{nmtfit} &
\begin{nmtfit}\begin{tikzpicture}[scale=0.33, thick]
    \useasboundingbox (-3.8,-3.4) rectangle (4.6,3.5);
    \draw[->] (-3.4,0) -- (4.2,0) node[below] {$x$};
    \draw[->] (0,-3.2) -- (0,3.0) node[left] {$y$};
    \node[below left] at (0,0) {\small $0$};
    \node[above left] at (3,0) {\small $3$};
    \node[below right] at (0,-2) {\small $-2$};
    \draw (-1.2,-2.8) -- (3.6,0.4);
\end{tikzpicture}\end{nmtfit} \\
\hline
\end{tabular}
\endgroup
\end{center}
\par\penalty-20\vspace{0.25cm}
\end{samepage}

\begin{samepage}
\noindent\zadnum Установіть відповідність між твердженням \mbox{(1--3)} та прямою, зображеною на рисунку \mbox{(А--Д)}, для якої це твердження є правильним. \nmtyear{2025}

\nmtnobreak
\nopagebreak\nbvspace{0.3cm}

\nmtnobreak
\noindent
\begin{minipage}[t]{0.60\textwidth}
\matchHead{Твердження}
    \matchItem{1}{не має спільних точок з функцією $y=\log_2(x-1)$}
\matchItem{2}{є графіком функції $y=\dfrac{2x}{3}-2$}
\matchItem{3}{кутовий коефіцієнт прямої є від'ємним числом}
\end{minipage}%
\hfill
\begin{minipage}[t]{0.30\textwidth}
    \nopagebreak\nbvspace{0.5cm}
    \matchingGrid
\end{minipage}

\nmtnobreak
\nopagebreak\nbvspace{0.5cm}
\matchHead{Пряма}

\nmtnobreak
% Таблиця з графіками
\begin{center}
\begingroup
\setlength{\tabcolsep}{3pt}
\begin{tabular}{|*{5}{c|}}
\hline
\textbf{А} & \textbf{Б} & \textbf{В} & \textbf{Г} & \textbf{Д} \\
\hline
\begin{nmtfit}\begin{tikzpicture}[scale=0.35]
    \draw[->] (-1.5,0) -- (3.5,0) node[below] {$x$};
    \draw[->] (0,-3) -- (0,2) node[left] {$y$};
    \node[below left] at (0,0) {\small $0$};
    \node[below] at (2,0) {\small $2$};
    \node[left] at (0,-2) {\small $-2$};
    \draw[thick] (-0.5,-2.5) -- (2.5,0.5); % y = x - 2
\end{tikzpicture}\end{nmtfit} &
\begin{nmtfit}\begin{tikzpicture}[scale=0.35]
    \draw[->] (-2.5,0) -- (2.5,0) node[below] {$x$};
    \draw[->] (0,-2) -- (0,3) node[left] {$y$};
    \node[below left] at (0,0) {\small $0$};
    \node[above right] at (0,1) {\small $1$};
    \draw (0.1,1) -- (0,1) -- (0.2,1.2) -- (0.2,1); % прямий кут
    \draw[thick] (-2.5,1) -- (2.5,1); % y = 1
\end{tikzpicture}\end{nmtfit} &
\begin{nmtfit}\begin{tikzpicture}[scale=0.35]
    \draw[->] (-2.5,0) -- (2.5,0) node[below] {$x$};
    \draw[->] (0,-2) -- (0,3) node[left] {$y$};
    \node[below right] at (0,0) {\small $0$};
    \node[below right] at (1,0) {\small $1$};
    \draw (1,0) -- (1.2,0) -- (1.2,0.2) -- (1,0.2); % прямий кут
    \draw[thick] (1,-2) -- (1,3); % x = 1
\end{tikzpicture}\end{nmtfit} &
\begin{nmtfit}\begin{tikzpicture}[scale=0.35]
    \draw[->] (-3.5,0) -- (2.5,0) node[below] {$x$};
    \draw[->] (0,-3) -- (0,2) node[left] {$y$};
    \node[below right] at (0,0) {\small $0$};
    \node[below] at (-2,0) {\small $-2$};
    \node[left] at (0,-2) {\small $-2$};
    \draw[thick] (-2.5,0.5) -- (0.5,-2.5); % y = -x - 2
\end{tikzpicture}\end{nmtfit} &
\begin{nmtfit}\begin{tikzpicture}[scale=0.35]
    \draw[->] (-1.5,0) -- (4,0) node[below] {$x$};
    \draw[->] (0,-3) -- (0,3) node[left] {$y$};
    \node[below left] at (0,0) {\small $0$};
    \node[below] at (3,0) {\small $3$};
    \node[left] at (0,-2) {\small $-2$};
    \draw[thick] (-1,-2.66) -- (4,0.66); % y = 2/3 x - 2
\end{tikzpicture}\end{nmtfit} \\
\hline
\end{tabular}
\endgroup
\end{center}
\par\penalty-20\vspace{0.25cm}
\end{samepage}

\clearpage
\sectionTitle{Відповіді}

\noindent{\small\color{gray!80!black}Для кожного завдання -- літери відповідей для пунктів 1, 2, 3. Відповіді обчислено й звірено з офіційними ключами там, де вони є.}\par\vspace{0.3cm}

\noindent\textbf{Частина 1}\par\vspace{0.1cm}

\begin{multicols}{5}\noindent
\mbox{1~--- ДБГ}\par
\mbox{2~--- АГБ}\par
\mbox{3~--- БАВ}\par
\mbox{4~--- ВАБ}\par
\mbox{5~--- АДБ}\par
\mbox{6~--- АДБ}\par
\mbox{7~--- ВБГ}\par
\mbox{8~--- АГБ}
\end{multicols}\par\vspace{0.2cm}

\noindent\textbf{Частина 2}\par\vspace{0.1cm}

\begin{multicols}{5}\noindent
\mbox{9~--- ГБА}\par
\mbox{10~--- АБВ}\par
\mbox{11~--- ГАД}\par
\mbox{12~--- БВА}\par
\mbox{13~--- БВА}\par
\mbox{14~--- АБГ}\par
\mbox{15~--- ВГА}\par
\mbox{16~--- АГВ}\par
\mbox{17~--- АГБ}\par
\mbox{18~--- ВГА}\par
\mbox{19~--- ДБА}\par
\mbox{20~--- ГБД}\par
\mbox{21~--- АГД}\par
\mbox{22~--- БАГ}\par
\mbox{23~--- ВБА}\par
\mbox{24~--- АБВ}\par
\mbox{25~--- ДАГ}\par
\mbox{26~--- ВБГ}\par
\mbox{27~--- АВБ}\par
\mbox{28~--- ГБА}\par
\mbox{29~--- БГД}\par
\mbox{30~--- БДГ}\par
\mbox{31~--- БГВ}\par
\mbox{32~--- ВДГ}\par
\mbox{33~--- ГВБ}\par
\mbox{34~--- БАД}\par
\mbox{35~--- ВГА}\par
\mbox{36~--- ГВА}\par
\mbox{37~--- ВГД}\par
\mbox{38~--- БАГ}\par
\mbox{39~--- ВБГ}\par
\mbox{40~--- ВАГ}\par
\mbox{41~--- БВД}\par
\mbox{42~--- АБВ}\par
\mbox{43~--- ВБД}\par
\mbox{44~--- ВБА}\par
\mbox{45~--- ГВА}\par
\mbox{46~--- ВАД}\par
\mbox{47~--- БВГ}\par
\mbox{48~--- ВБГ}\par
\mbox{49~--- АБВ}\par
\mbox{50~--- ДБА}\par
\mbox{51~--- ВАБ}\par
\mbox{52~--- ВАД}\par
\mbox{53~--- ВБА}\par
\mbox{54~--- БВД}\par
\mbox{55~--- ГАБ}\par
\mbox{56~--- ВГБ}\par
\mbox{57~--- АГВ}\par
\mbox{58~--- АГБ}\par
\mbox{59~--- ВБА}\par
\mbox{60~--- ДГБ}\par
\mbox{61~--- ВБД}\par
\mbox{62~--- БГВ}\par
\mbox{63~--- ДАВ}\par
\mbox{64~--- ДАБ}\par
\mbox{65~--- ВДГ}\par
\mbox{66~--- БГД}\par
\mbox{67~--- ГАД}\par
\mbox{68~--- ГАБ}\par
\mbox{69~--- ВДГ}
\end{multicols}\par\vspace{0.2cm}

\end{document}
